\subsection{INNER PRODUCTS: CASE OF COGEBRAS}

Let \(E = \bigoplus_{p \geq 0} E_p\) be a coassociative counital graded cogebra. Then we know (no. 2, Propositions 1 and 3) that the graded dual \(E^{*gr}\) has (with the convention on graduations made at the beginning of the paragraph) a graded algebra structure of type N over A, the product of two elements \(u, v\) of this algebra being defined by \(uv = m \circ (u \otimes v) \circ c\) , where \(c: E \to E \otimes_A E\) is the coproduct and \(m: A \otimes_A A \to A\) defines the multiplication. In other words, if, for \(x \in E\) , \(c(x) = \sum_i y_i \otimes z_i\) , we can write (canonically identifying \(A \otimes_A E\) and \(E\) )

\[
\begin{array}{r l} \langle x, u v \rangle = (u v) (x) & = \sum_ {i} u (y _ {i}) v (z _ {i}) = v \left(\sum_ {i} u (y _ {i}) z _ {i}\right) \\ & = v (((u \otimes 1 _ {\mathbb {E}}) \circ c) (x)) = \langle ((u \otimes 1 _ {\mathbb {E}}) \circ c) (x), v \rangle . \end{array}
\]

This can be interpreted by saying that, for all u homogeneous of degree p in \(E^{*gr}\) , the left multiplication \(e(u): v \mapsto uv\) in \(E^{*gr}\) is the graded transpose of the graded endomorphism of degree -p

\[
i (u) = (u \otimes 1 _ {\mathrm{E}}) \circ c\tag{46}
\]

of E; hence, in the above notation,

\[
(i (u)) (x) = \sum_ {i} u (y _ {i}) z _ {i}.
\]

Formula (46) also defines an element \(i(u)\in\mathrm{Endgr}_{\mathbf{A}}(\mathbf{E})\) for every element \(u\in\mathbf{E}^{*\mathbf{g}\mathbf{r}}\) ; for all \(x\in\mathbf{E}\) and all \(u\in\mathbf{E}^{*\mathbf{g}\mathbf{r}}\) , we write

\[
x \llcorner u = (i (u)) (x)\tag{47}
\]

so that, for u and v in E*gr,

\[
\langle x, u v \rangle = \langle x \llcorner u, v \rangle .
\]

The element \(x \llcorner u\) of E is called the right inner product of x by u.

Here again the element u which ``operates'' on x is placed at the free end of the horizontal line in L.

For any two elements \(u, v\) of \(\mathbf{E}^{*\mathrm{gr}}\) ,

\[
x \llcorner (u v) = (x \llcorner u) \llcorner v,\tag{48}
\]

in other words

\[
i (u v) = i (v) \circ i (u).
\]

As above let \(c(x) = \sum_{i} y_{i} \otimes z_{i}\) , so that \(x \llcorner (uv) = \sum_{i} (uv)(y_{i})z_{i}\) . If

\[
c (y _ {i}) = \sum_ {j} y _ {i j} ^ {\prime} \otimes y _ {i j} ^ {\prime \prime},
\]

then

\[
x \llcorner (u v) = \sum_ {i, j} u \left(y _ {i j} ^ {\prime}\right) v \left(y _ {i j} ^ {\prime \prime}\right) z _ {i}.\tag{49}
\]

On the other hand, if \(c(z_{i}) = \sum_{k} z_{ik}' \otimes z_{ik}''\) , then

\[
(x \llcorner u) \llcorner v = \sum_ {i, k} u (y _ {i}) v (z _ {i k} ^ {\prime}) z _ {i k} ^ {\prime \prime}.\tag{50}
\]

Now, the coassociativity of E shows that (no. 2, Proposition 1)

\[
\sum_ {i, j} y _ {i j} ^ {\prime} \otimes y _ {i j} ^ {\prime \prime} \otimes z _ {i} = \sum_ {i, k} y _ {i} \otimes z _ {i k} ^ {\prime} \otimes z _ {i k} ^ {\prime \prime}\tag{51}
\]

and the equality of expressions (49) and (50) follows from the fact that they are respectively the image of the left and right hand sides of (51) under the linear mapping f from \(E \otimes E \otimes E\) to E such that \(f(x \otimes y \otimes z) = u(x)v(y)z\) .

We recall on the other hand (no. 2, Proposition 3) that the unit element of the algebra \(\mathbf{E}^{*\mathrm{gr}}\) is the linear form \(e: x \mapsto \gamma(x).1\) ; hence

\[
x \llcorner e = \sum_ {i} \gamma (y _ {i}) z _ {i} = x
\]

by virtue of the definition of counit. As the mapping \(u \mapsto i(u)\) is linear, it is seen that the external law of composition \((u, x) \mapsto x \llcorner u\) defines a right E*gr-module structure on E.

Similarly we define, for all \(u \in \mathrm{E}^{*\mathbf{g}\mathbf{r}}\) , the endomorphism of \(\mathbf{E}\)

\[
i ^ {\prime} (u) = \left(1 _ {E} \otimes u\right) \circ c\tag{52}
\]

and, for all \(x \in \mathbf{E}\) , we write

\[
\left(i ^ {\prime} (u)\right) (x) = u \lrcorner x\tag{53}
\]

and this element of E is called the left inner product of x by u. As above it is seen that the external law \((u, v) \mapsto u \lrcorner x\) defines a left \(\mathrm{E}^{*gr}\) -module structure on E. Moreover, these two structures are compatible, in other words,

\[
(u \lrcorner x) \llcorner v = u \lrcorner (x \llcorner v)\tag{54}
\]

for \(u, v\) in E*gr (II, § 1, no. 14). With the same notation as above, the left hand side of (54) is \(\sum_{i,j} u(z_i) v(y_{ij}') y_{ij}''\) and the right hand side is \(\sum_{i,k} v(y_i) u(z_{ik}') z_{ik}''\) ; their equality follows from the fact that they are the respective images of the left and right hand sides of (51) under the linear mapping \(g\) of E ⊗ E ⊗ E into E such that \(g(x \otimes y \otimes z) = v(x) u(z)y\) .

It is therefore seen that the two external laws of composition on E defines on this set an \((\mathrm{E}^{*}\mathrm{gr}, \mathrm{E}^{*}\mathrm{gr})\) -bimodule structure.

When the cogebra E is cocommutative, then \(u \lrcorner x = x \llcorner u\) for all \(x \in \mathbf{E}\) and \(u \in \mathrm{E}^{*\mathrm{gr}}\) ; when it is anticocommutative (§ 4, no. 9) and \(u \in \mathrm{E}_p^*\) and \(x \in \mathrm{E}_q\) , we can write \(c(x) = \sum_{0 \leqslant j \leqslant q} \left( \sum_i y_{ij} \otimes z_{i,q-j} \right)\) with \(y_{ij}\) and \(z_{ij}\) in \(\mathrm{E}_j\) for all \(j\) and then by hypothesis

\[
\sum_ {i} z _ {i j} \otimes y _ {i, q - j} = (- 1) ^ {j (q - j)} \sum_ {i} y _ {i j} \otimes z _ {i, q - j}.
\]

By definition, \(x \llcorner u = \sum_{0 \leqslant j \leqslant q} \left( \sum_{i} u(y_{ij}) z_{i,q-j} \right)\) and

\[
u \lrcorner x = \sum_ {0 \leqslant j \leqslant q} \left(\sum_ {i} u (z _ {i, q - j}) y _ {i j}\right).
\]

As \(u(y_{ij}) = 0\) (resp. \(u(z_{i,q-j}) = 0\) unless \(j = p\) (resp. \(q - j = p\) ), it is seen by the above that \(u \lrcorner x = (-1)^{p(q-p)}x \llcorner u\) .

Finally, let \(\phi: \mathrm{E} \to \mathrm{F}\) be a graded cogebra morphism; then it has been seen (no. 2) that the graded transpose \(^t\phi: \mathrm{F}^{*\mathrm{gr}} \to \mathrm{E}^{*\mathrm{gr}}\) is a graded algebra homomorphism; therefore, for \(x \in \mathrm{E}, u, v\) in \(\mathrm{F}^{*\mathrm{gr}}\) ,

\[
\begin{array}{r l} \langle \phi (x \llcorner {} ^ {t} \phi (u)), v \rangle & = \langle x \llcorner {} ^ {t} \phi (u), ^ {t} \phi (v) \rangle = \langle x, ^ {t} \phi (u) ^ {t} \phi (v) \rangle = \langle x, ^ {t} \phi (u v) \rangle \\ & = \langle \phi (x), u v \rangle = \langle \phi (x) \llcorner u, v \rangle \end{array}
\]

whence

\[
\phi (x) \llcorner u = \phi (x \llcorner {} ^ {t} \phi (u));\tag{55}
\]

and similarly

\[
u \lrcorner \phi (x) = \phi (^ {t} \phi (u) \lrcorner x).\tag{56}
\]

In other words, \(\phi\) is an \(\mathbf{F}^{*\mathrm{gr}}\) -homomorphism of the \((\mathrm{E}^{*\mathrm{gr}}, \mathrm{E}^{*\mathrm{gr}})\) -bimodule E into the \((\mathrm{F}^{*\mathrm{gr}}, \mathrm{F}^{*\mathrm{gr}})\) -bimodule F, relative to the ring homomorphism

\[
{ } ^ { t } \phi : \mathrm{F} ^ { * \mathrm{gr} } \rightarrow \mathrm{E} ^ { * \mathrm{gr} } .
\]

Examples. In particular the above can be applied when E is one of the graded cogebras \(\mathsf{T}(\mathbf{M})\) , \(\mathsf{S}(\mathbf{M})\) or \(\bigwedge(\mathbf{M})\) for an A-module M (no. 1, Examples 5, 6 and 7). To find explicitly the bimodule structures thus obtained, we again identify a homogeneous element f of degree n in \(\mathsf{T}(\mathbf{M})^{*\mathbf{g}\mathbf{r}}\) (resp. \(\mathsf{S}(\mathbf{M})^{*\mathbf{g}\mathbf{r}}\) , resp. \(\bigwedge(\mathbf{M})^{*\mathbf{g}\mathbf{r}}\) ) with an n-linear form (resp. symmetric n-linear form, resp. alternating n-linear form, also called an n-form) on \(M^{n}\) . It suffices to express the products

\[
\left(x _ {1} \otimes x _ {2} \otimes \dots \otimes x _ {p}\right) \llcorner f (\text { resp. } (x _ {1} x _ {2} \dots x _ {p}) \llcorner f,
\]

resp. \((x_{1} \wedge x_{2} \wedge \cdots \wedge x_{p}) \llcorner f)\) for every finite sequence \((x_{i})_{1 \leqslant i \leqslant p}\) of elements of M and the analogues for the left inner product. Now, definitions (46) and (52) and formulae (3), (6) and (9) of no. 1 give respectively:

\[
\left\{ \begin{array}{l} (x _ {1} \otimes x _ {2} \otimes \dots \otimes x _ {p}) \llcorner f = f \lrcorner (x _ {1} \otimes x _ {2} \otimes \dots \otimes x _ {p}) = 0 \\ (x _ {1} x _ {2} \ldots x _ {p}) \llcorner f = f \lrcorner (x _ {1} x _ {2} \ldots x _ {p}) = 0 \quad \text {for} \quad p <   n. \\ (x _ {1} \wedge x _ {2} \wedge \dots \wedge x _ {p}) \llcorner f = f \lrcorner (x _ {1} \wedge x _ {2} \wedge \dots \wedge x _ {p}) = 0 \end{array} \right.\tag{57}
\]

For \(p \geqslant n\) , we have respectively

\begin{enumerate}
\def\labelenumi{(\arabic{enumi})}
\setcounter{enumi}{57}
\tightlist
\item
\end{enumerate}

\[
\left(x _ {1} \otimes x _ {2} \otimes \dots \otimes x _ {p}\right) \llcorner f = f \left(x _ {1}, \dots , x _ {n}\right) x _ {n + 1} \otimes \dots \otimes x _ {p}\tag{59}
\]

\[
\left(x _ {1} x _ {2} \dots x _ {p}\right) \llcorner f = \sum_ {\sigma} f \left(x _ {\sigma (1)}, \dots , x _ {\sigma (n)}\right) x _ {\sigma (n + 1)} \dots x _ {\sigma (p)}\tag{60}
\]

\[
\left(x _ {1} \wedge x _ {2} \wedge \dots \wedge x _ {p}\right) \llcorner f = \sum_ {\sigma} \varepsilon_ {\sigma} f \left(x _ {\sigma (1)}, \dots , x _ {\sigma (n)}\right) x _ {\sigma (n + 1)} \wedge \dots \wedge x _ {\sigma (p)}
\]

(where, in (59) and (60), the summations are taken over the permutations \(\sigma \in \mathfrak{S}_p\) which are increasing on each of the intervals \([1, n]\) and \([n + 1, p]\) of \(\mathbf{N}\) ); and similarly

\begin{enumerate}
\def\labelenumi{(\arabic{enumi})}
\setcounter{enumi}{60}
\tightlist
\item
\end{enumerate}

\[
f \lrcorner (x _ {1} \otimes x _ {2} \otimes \dots \otimes x _ {p}) = f (x _ {p - n + 1}, \dots , x _ {p}) x _ {1} \otimes x _ {2} \otimes \dots \otimes x _ {p - n}\tag{62}
\]

\[
f \lrcorner (x _ {1} x _ {2} \dots x _ {p}) = \sum_ {\sigma} f (x _ {\sigma (p - n + 1)}, \dots , x _ {\sigma (p)}) x _ {\sigma (1)} \dots x _ {\sigma (p - n)}\tag{63}
\]

\[
f \lrcorner (x _ {1} \wedge x _ {2} \wedge \dots \wedge x _ {p}) =
\]

\[
\sum_ {\sigma} \varepsilon_ {\sigma} f (x _ {\sigma (p - n + 1)}, \dots , x _ {\sigma (p)}) x _ {\sigma (1)} \wedge \dots \wedge x _ {\sigma (p - n)}
\]

(where, in (62) and (63), the summations are taken over the permutations \(\sigma \in \mathfrak{S}_p\) which are increasing on each of the intervals \([1, p - n]\) and \([p - n + 1, p]\) of \(\mathbf{N}\) ).

\subsection{INNER PRODUCTS: CASE OF BIGEBRAS}

Let E be a graded bigebra (resp. skew graded bigebra) (no. 4, Definition 3); then the results of nos. 6 and 7 can be applied to define the right (resp. left) inner products \(x \llcorner u \in \mathbf{E}\) and \(u \llcorner x \in \mathbf{E}^{*\mathbf{gr}}\) (resp. \(u \lrcorner x \in \mathbf{E}\) and \(x \lrcorner u \in \mathbf{E}^{*\mathbf{gr}}\) )

for all \(x \in \mathbf{E}\) and all \(u \in \mathbf{E}^{*\mathbf{gr}}\) . Thus an (E, E)-bimodule structure and an \((\mathbf{E}^{*\mathbf{gr}}, \mathbf{E}^{*\mathbf{gr}})\) -bimodule structure are obtained on E. Further:

PROPOSITION 9. Let E be a graded bigebra (resp. skew graded bigebra). For every element x of degree 1 in E, the left and right inner products by x are derivations (resp. antiderivations) (§ 10, no. 2) of the algebra E*gr.

In the notation of no. 6, for every homogeneous element \(x\) of degree 1 in a graded bigebra (resp. a skew graded bigebra) E,

\[
c (x) = x \otimes 1 + 1 \otimes x,
\]

by Proposition 5 of no. 3 and the fact that \(c\) is a homomorphism of degree 0. Suppose first that \(E\) is a graded bigebra. For all \(y \in E\) , by definition

\[
\langle (u v) \llcorner x, y \rangle = \langle u v, x y \rangle = m ((u \otimes v) (c (x y)))
\]

and since \(c\) is an algebra homomorphism, \(c(xy) = c(x)c(y)\) . Let \(c(y) = \sum_{i} s_i \otimes t_i\) with \(s_i\) and \(t_i\) in \(E\) ; therefore

\[
c (x y) = \sum_ {i} (x s _ {i}) \otimes t _ {i} + \sum_ {i} s _ {i} \otimes (x t _ {i}).
\]

Hence \(\langle (uv) \llcorner x, y \rangle = \sum_{i} u(xs_i)v(t_i) + \sum_{i} u(s_i)v(xt_i)\) . But we may write

\[
\sum_ {i} u (x s _ {i}) v (t _ {i}) = m (((u \llcorner x) \otimes v) (c (y))) = \langle (u \llcorner x) v, y \rangle
\]

and similarly

\[
\sum_ {i} u (s _ {i}) v (x t _ {i}) = m ((u \otimes (v \llcorner x)) (c (y))) = \langle u (v \llcorner x), y \rangle ,
\]

whence, returning to the notation \(i(x)\) for the inner product,

\[
(i (x)) (u v) = ((i (x)) (u)) v + u ((i (x)) (v))\tag{64}
\]

which proves that \(i(x)\) is a derivation in \(\mathbf{E}^{*}\mathbf{gr}\) .

Suppose now that \(\mathbf{E}\) is a skew graded bigebra, that \(u\in \mathbf{E}_p^*\) , \(v\in \mathbf{E}_q^*\) and \(y\in \mathbf{E}_r\) ; then we can write

\[
c (y) = \sum_ {0 \leqslant j \leqslant r} \left(\sum_ {i} s _ {i j} \otimes t _ {i, r - j}\right)
\]

where the \(s_{ij}\) and \(t_{ij}\) belong to \(\mathbf{E}_j\) ; by definition of the product in \(\mathbf{E}^{\mathfrak{g}} \otimes_{\mathbf{A}} \mathbf{E}\) , then

\[
c (x y) = c (x) c (y) = \sum_ {0 \leqslant j \leqslant r} \left(\sum_ {i} \left(x s _ {i j}\right) \otimes t _ {i, r - j} + (- 1) ^ {j} \sum_ {i} s _ {i j} \otimes \left(x t _ {i, r - j}\right)\right)
\]

whence this time

\[
\langle (u v) \llcorner x, y \rangle = \sum_ {0 \leqslant j \leqslant r} \left(\sum_ {i} u \left(x s _ {i j}\right) v \left(t _ {i, r - j}\right) + (- 1) ^ {j} \sum_ {i} u \left(s _ {i j}\right) v \left(x t _ {i, r - j}\right)\right).
\]

Then also \(\sum_{0\leqslant j\leqslant r}\left(\sum_{i}u(xs_{ij})v(t_{i,r - j})\right) = \langle (u\llcorner x)v,y\rangle\) . On the other hand, \(u(s_{ij}) = 0\) unless \(j = -p\) and hence we may also write

\[
\sum_ {0 \leqslant j \leqslant r} (- 1) ^ {j} \left(\sum_ {i} u (s _ {i j}) v (x t _ {i, r - j})\right) = (- 1) ^ {p} \langle u (v \llcorner x), y \rangle .
\]

We therefore conclude that

\[
(i (x)) (u v) = ((i (x)) (u)) v + (- 1) ^ {p} u ((i (x)) (v)),\tag{65}
\]

in other words \(i(x)\) is an antiderivation in \(\mathrm{E}^{*\mathrm{gr}}\) . The assertions relating to the left inner product by an element \(x\) of degree 1 in E are proved similarly.

Remarks. (1) Let E be a graded bigebra over A and N, N', N'' three graded A-modules. Let m be an A-bilinear mapping of N × N' into N''; for u \ensuremath{\in} Homgr\_A(E, N) and v \ensuremath{\in} Homgr\_A(E, N'), let u.v denote the graded homomorphism m ∘ (u ⊗ v) ∘ c of E into N''. On the other hand, let i(x) denote the (right or left) inner product by x \ensuremath{\in} E in the A-modules Homgr\_A(E, N), Homgr\_A(E, N') and Homgr\_A(E, N'\,'). Then, if x is of degree 1,

\[
(i (x)) (u. v) = ((i (x)) (u)). v + u. ((i (x)) (v))
\]

for all \(u \in \operatorname{Homgr}_{\mathbf{A}}(\mathbf{E}, \mathbf{N})\) and \(v \in \operatorname{Homgr}_{\mathbf{A}}(\mathbf{E}, \mathbf{N}')\) .

Under the same conditions, if E is a skew graded bigebra and u is homogeneous of degree p, then

\[
(i (x)) (u. v) = ((i (x)) (u)). v + (- 1) ^ {p} u. ((i (x)) (v)).
\]

The proofs are the same as in Proposition 9.

\begin{enumerate}
\def\labelenumi{(\arabic{enumi})}
\setcounter{enumi}{1}
\tightlist
\item
  The same argument as in the above proof proves, more generally, that for all \(x \in \mathbf{E}\) , if \(c(x) = \sum_{j} x_j' \otimes x_{j'}'\) then, for all \(u, v\) in \(\mathbf{E}^{*\mathrm{gr}}\) , the ``Leibniz formula''
\end{enumerate}

\[
(i (x)) (u v) = \sum_ {j} (i (x _ {j} ^ {\prime})) (u) \cdot (i (x _ {j} ^ {\prime \prime})) (v)
\]

holds. In particular, for every primitive element of a graded bigebra E, that is such that \(c(x) = x \otimes 1 + 1 \otimes x\) , \(i(x)\) is a derivation of \(\mathrm{E}^{*\mathrm{gr}}\) .

PROPOSITION 10. Let E be a graded bigebra (resp. skew graded bigebra). For every element f of degree 1 in E*gr, the left and right inner products are derivations (resp. antiderivations) of the algebra E.

Let \(x \in \mathrm{E}_p, y \in \mathrm{E}_p\) ( \(p \geqslant 1, q \geqslant 1\) ). By Proposition 5 of no. 3, we can write

\[
c (x) = x \otimes 1 + \sum_ {1 \leqslant j \leqslant p - 1} \left(\sum_ {i} x _ {i j} ^ {\prime} \otimes x _ {i j, p - j} ^ {\prime \prime}\right) + 1 \otimes x
\]

\[
c (y) = y \otimes 1 + \sum_ {1 \leqslant k \leqslant q - 1} \left(\sum_ {i} y _ {i k} ^ {\prime} \otimes y _ {i, q - k} ^ {\prime \prime}\right) + 1 \otimes y
\]

where \(x_{ij}'\) and \(x_{ij}''\) belong to \(\mathrm{E}_j\) , \(y_{ik}'\) and \(y_{ik}''\) to \(\mathrm{E}_k\) . If \(\mathrm{E}\) is a graded bigebra, the component of \(c(xy) = c(x)c(y)\) belonging to \(\mathrm{E}_1 \otimes \mathrm{E}\) , is equal to

\[
\sum_ {i} x _ {i, 1} ^ {\prime} \otimes x _ {i, p - 1} ^ {\prime \prime} y + \sum_ {i} y _ {i, 1} ^ {\prime} \otimes x y _ {i, q - 1} ^ {\prime \prime}
\]

and hence by definition

\[
\begin{array}{r l} (x y) \llcorner f & = \sum_ {i} f (x _ {i, 1} ^ {\prime}) x _ {i, p - 1} ^ {\prime \prime} y + \sum_ {i} f (y _ {i, 1} ^ {\prime}) x y _ {i, q - 1} ^ {\prime \prime} \\ & = (x \llcorner f) y + x (y \llcorner f) \end{array}
\]

and the right inner product by \(f\) is a derivation. If on the other hand \(E\) is a skew graded bigebra, the component of \(c(xy)\) belonging to \(E_1 \otimes E\) is equal to

\[
\sum_ {i} x _ {i, 1} ^ {\prime} \otimes x _ {i, p - 1} ^ {\prime \prime} y + (- 1) ^ {p} \sum_ {i} y _ {i, 1} ^ {\prime} \otimes x y _ {i, q - 1} ^ {\prime \prime}
\]

and this time we obtain

\[
(x y) \llcorner f = (x \llcorner f) y + (- 1) ^ {p} x (y \llcorner f)
\]

which shows that \(i(f)\) is then an antiderivation. The argument is similar for the left inner product by \(f\) .

Examples. Propositions 9 and 10 apply in particular to the graded bigebra \(\mathsf{S}(\mathbf{M})\) and the skew graded bigebra \(\bigwedge (\mathbf{M})\) . The inner products by elements of degree 1 in \(\mathsf{S}(\mathbf{M})\) (resp. \(\mathsf{S}(\mathbf{M})^{*\mathrm{gr}}\) ) are derivations which commute with one another, since \(\mathsf{S}(\mathbf{M})\) (resp. \(\mathsf{S}(\mathbf{M})^{*\mathrm{gr}}\) ) is commutative.

Similarly, the inner products by elements of degree 1 in \(\bigwedge (\mathbf{M})\) (resp. \(\bigwedge (\mathbf{M})^{*\mathrm{gr}}\) ) are antiderivations, which are of zero square, for the square of an element of degree 1 in the algebra \(\bigwedge (\mathbf{M})\) (resp. \(\bigwedge (\mathbf{M})^{*\mathrm{gr}}\) ) is zero.

\subsection{\texorpdfstring{INNER PRODUCTS BETWEEN \(\mathsf{T}(\mathbf{M})\) AND \(\mathsf{T}(\mathbf{M}^{*})\), \(\mathsf{S}(\mathbf{M})\) AND \(\mathsf{S}(\mathbf{M}^{*})\), \(\bigwedge(\mathbf{M})\) AND \(\bigwedge(\mathbf{M}^{*})\)}{Inner products between T(M) and T(M*), S(M) and S(M*), wedge(M) and wedge(M*)}}

The right inner product defines on \(\mathsf{T}(\mathbf{M})\) (resp. \(\mathsf{S}(\mathbf{M})\) , resp. \(\bigwedge (\mathbf{M})\) ) a right module structure over the algebra \(\mathsf{T}(\mathbf{M})^{*\mathfrak{g}\mathfrak{r}}\) (resp. \(\mathsf{S}(\mathbf{M})^{*\mathfrak{g}\mathfrak{r}}\) , resp. \(\bigwedge (\mathbf{M})^{*\mathfrak{g}\mathfrak{r}}\) ) (no. 7, Examples). Using the canonical homomorphisms \(\theta_{\mathsf{T}}\) (resp. \(\theta_{\mathsf{S}}\) , resp. \(\theta_{\wedge}\) ) of no. 5, we derive

a right \(\mathsf{T}(\mathbf{M}^{*})\) -module structure on \(\mathsf{T}(\mathbf{M})\)

a right \(\mathbf{S}(\mathbf{M}^{*})\) -module structure on \(\mathbf{S}(\mathbf{M})\)

a left \(\bigwedge (\mathbf{M}^{*})\) -module structure on \(\bigwedge (\mathbf{M})\) .

The external law of any of these structures is also denoted by

\[
(z ^ {*}, t) \mapsto i (z ^ {*}). t
\]

(by an abuse of language); we also write \(t \llcorner z^*\) instead of \(i(z^*)\) . \(t\) in the case of \(\mathsf{T}(\mathbf{M})\) or \(\mathsf{S}(\mathbf{M})\) ; on the other hand, we write \(z^* \lrcorner t\) in the case of \(\bigwedge (\mathbf{M})\) and say that this is a left inner product of \(t\) by \(z^*\) , since then we have a left \(\bigwedge (\mathbf{M}^*)\) -module law. For \(z^*\) homogeneous of degree \(n\) and \(t\) homogeneous of degree \(p\) , \(i(z^*)\) . \(t = 0\) if \(p < n\) and, for \(x_i \in \mathbf{M} (1 \leqslant i \leqslant p)\) , \(x_j^* \in \mathbf{M}^* (1 \leqslant j \leqslant n)\) and \(p \geqslant n\) , by virtue of formulae (58), (59) and (60) of no. 7,

\[
\begin{array}{r l} i (x _ {1} ^ {*} \otimes x _ {2} ^ {*} \otimes \dots \otimes x _ {n} ^ {*}) \cdot (x _ {1} \otimes x _ {2} \otimes \dots \otimes x _ {p}) \\ & = \left(\prod_ {j = 1} ^ {n} \langle x _ {j} ^ {*}, x _ {j} \rangle\right) x _ {n + 1} \otimes \dots \otimes x _ {p} \end{array}\tag{66}
\]

\[
i \left(x _ {1} ^ {*} x _ {2} ^ {*} \dots x _ {n} ^ {*}\right). (x _ {1} x _ {2} \dots x _ {p}) = \sum_ {\sigma} \left(\prod_ {j = 1} ^ {n} \langle x _ {j} ^ {*}, x _ {\sigma (j)} \rangle\right) x _ {\sigma (n + 1)} \dots x _ {\sigma (p)}\tag{67}
\]

\[
\begin{array}{l} i \left(x _ {1} ^ {*} \wedge x _ {2} ^ {*} \wedge \dots \wedge x _ {n} ^ {*}\right). \left(x _ {1} \wedge x _ {2} \wedge \dots \wedge x _ {p}\right) \\ = (- 1) ^ {n (n - 1) / 2} \sum_ {\sigma} \varepsilon_ {\sigma} \left(\prod_ {j = 1} ^ {n} \langle x _ {j} ^ {*}, x _ {\sigma (j)} \rangle\right) x _ {\sigma (n + 1)} \wedge \dots \wedge x _ {\sigma (p)} \end{array}\tag{68}
\]

where, the formulae (67) and (68), \(\sigma\) runs through the set of permutations \(\sigma \in \mathfrak{S}_p\) which are increasing on the intervals \([1, n]\) and \([n + 1, p]\) . We can also write, in the inner product notation,

\[
\begin{array}{l l} \langle t \llcorner u ^ {*}, v ^ {*} \rangle = \langle t, \theta_ {\mathsf {T}} (u ^ {*} v ^ {*}) \rangle & \text {for} \quad t \in \mathsf {T} (\mathbf {M}), u ^ {*}, v ^ {*} \quad \text {in} \quad \mathsf {T} (\mathbf {M} ^ {*}) \\ \langle t \llcorner u ^ {*}, v ^ {*} \rangle = \langle t, \theta_ {\mathsf {S}} (u ^ {*} v ^ {*}) \rangle & \text {for} \quad t \in \mathsf {S} (\mathbf {M}), u ^ {*}, v ^ {*} \quad \text {in} \quad \mathsf {S} (\mathbf {M} ^ {*}) \\ \langle v ^ {*}, u ^ {*} \lrcorner t \rangle = \langle \theta_ {\wedge} (u ^ {*} \wedge v ^ {*}), t \rangle & \text {for} \quad t \in \bigwedge (\mathbf {M}), u ^ {*}, v ^ {*} \quad \text {in} \quad \bigwedge (\mathbf {M} ^ {*}). \end{array}
\]

We leave to the reader the task of finding explicitly the analogous formulae for left inner products, this time using formulae (61), (62) and (63).

The above can be applied with M replaced by its dual \(M^{*}\) ; \(M^{*}\) must then be replaced by the bidual \(M^{**}\) and \(T(M^{*})\) , for example, thus has a right module structure over the algebra \(T(M^{**})\) . But the canonical mapping \(c_{M}: M \to M^{**}\) defines an algebra homomorphism \(T(c_{M}): T(M) \to T(M^{**})\) , by means of which \(T(M^{*})\) has a right \(T(M)\) -module structure. Similarly \(S(M^{*})\) (resp. \(\bigwedge(M^{*})\) ) has a right \(S(M)\) -module (resp. left \(\bigwedge(M)\) -module) structure. The explicit formulae giving the external laws of these modules are derived immediately from the above by exchanging the roles of M and \(M^{*}\) . Note that, for all \(x \in M\) , \(i(x)\) is always a derivation (resp. antiderivation of zero square) of the graded algebra \(S(M^{*})\) (resp. \(\bigwedge(M^{*})\) ).

PROPOSITION 11. The canonical homomorphism \(\theta_{\mathsf{T}}: \mathsf{T}(\mathbf{M}^{*}) \to \mathsf{T}(\mathbf{M})^{*\mathsf{gr}}\) (resp. \(\theta_{\mathsf{S}}: \mathsf{S}(\mathbf{M}) \to \mathsf{S}(\mathbf{M})^{*\mathsf{gr}}\) , resp. \(\theta_{\wedge}: \bigwedge (\mathbf{M}^{*}) \to \bigwedge (\mathbf{M})^{*\mathsf{gr}}\) ) is a right \(\mathsf{T}(\mathbf{M})\) -module (resp. right \(\mathsf{S}(\mathbf{M})\) -module, resp. left \(\bigwedge (\mathbf{M})\) -module) homomorphism.

We show first that, for \(z^{*} \in \mathsf{T}(\mathbf{M}^{*})\) and \(t \in \mathsf{T}(\mathbf{M})\) ,

\[
\theta_ {T} (z ^ {*} \llcorner t) = \theta_ {T} (z ^ {*}) \llcorner t.\tag{69}
\]

Since M is a generating system of the algebra \(\mathsf{T}(\mathbf{M})\) , we need only prove (69) when \(t = x \in M\) ; moreover we can restrict our attention to the case where \(z^{*} = x_{1}^{*} \otimes x_{2}^{*} \otimes \cdots \otimes x_{p}^{*}\) , where the \(x_{j}^{*} \in M^{*}\) , and then, by (66) with the roles of M and \(M^{*}\) interchanged, \(z^{*} \llcorner x = \langle x, x_{1}^{*} \rangle x_{2}^{*} \otimes \cdots \otimes x_{p}^{*}\) . Therefore, for all \(y_{2}, \ldots, y_{p}\) in M,

\[
\begin{array}{r l} \langle \theta_ {T} (z ^ {*} \llcorner x), y _ {2} \otimes y _ {3} \otimes \dots \otimes y _ {p} \rangle & = \langle x, x _ {1} ^ {*} \rangle \prod_ {j = 2} ^ {p} \langle y _ {j}, x _ {j} ^ {*} \rangle \\ & = \langle \theta_ {T} (z ^ {*}), x \otimes y _ {2} \otimes \dots \otimes y _ {p} \rangle \\ & = \langle \theta_ {T} (z ^ {*}) \llcorner x, y _ {2} \otimes \dots \otimes y _ {p} \rangle \end{array}
\]

whence (69).

We prove secondly that for \(z^{*} \in \mathbf{S}(\mathbf{M}^{*})\) and \(t \in \mathbf{S}(\mathbf{M})\) ,

\[
\theta_ {S} (z ^ {*} \llcorner t) = \theta_ {S} (z ^ {*}) \llcorner t.\tag{70}
\]

As above, we can limit ourselves to the case \(t = x \in \mathbf{M}\) . But further, here \(i(x)\) is a derivation of \(\mathbf{S}(\mathbf{M}^*)\) and a derivation of \(\mathbf{S}(\mathbf{M})^{*\text{gr}}\) . Therefore (§ 10, no. 7, Corollary to Proposition 9) it suffices to verify (70) for \(z^* = x^* \in \mathbf{M}^*\) , since \(\mathbf{M}^*\) is a generating system of \(\mathbf{S}(\mathbf{M}^*)\) ; but this is trivial, the two sides then being equal to \(\langle x^*, x \rangle\) . A similar argument proves the relation

\[
\theta_ {\wedge} (t \lrcorner z ^ {*}) = t \lrcorner \theta_ {\wedge} (z ^ {*})\tag{71}
\]

for \(z^{*} \in \bigwedge(\mathbf{M}^{*})\) and \(t \in \bigwedge(\mathbf{M})\) : observe then that, for \(x \in \mathbf{M}\) , \(i(x)\) is an antiderivation in \(\bigwedge(\mathbf{M}^{*})\) as well as in \(\bigwedge(\mathbf{M})^{*\text{er}}\) and use § 10, no. 7, Corollary to Proposition 9. There is an analogous result for left inner products.

\subsection{EXPLICIT FORM OF INNER PRODUCTS IN THE CASE OF A FINITELY GENERATED FREE MODULE}

Let \(\mathbf{M}\) be a finitely generated free A-module, \((e_i)_{1\leq i\leq n}\) a basis of \(\mathbf{M}\) and \((e_i^*)_{1\leq i\leq n}\) the dual basis of \(\mathbf{M}^*\) . For every finite sequence \(s = (i_1,\ldots ,i_p)\) of elements of \([1,n]\) , let \(e_s = e_{i_1}\otimes e_{i_2}\otimes \dots \otimes e_{i_p}\) (resp. \(e_s^* = e_{i_1}^*\otimes \dots \otimes e_{i_p}^*\) ). We know (§ 5, no. 5, Theorem 1) that the \(e_s\) form a basis of the A-module \(\mathsf{T}(\mathbf{M})\) and the \(e_s^*\) a basis of the A-module \(\mathsf{T}(\mathbf{M}^*)\) . If \(s,t\) are two finite sequences of elements of \([1,n]\) , let \(s.t\) denote the sequence obtained as follows: if \(s = (i_1,\ldots ,i_p)\) and \(t = (j_1,\ldots ,j_q)\) , \(s.t\) is the sequence \((i_1,\ldots ,i_p,j_1,\ldots ,j_q)\) with \(p + q\) terms. Then \(e_{s,t} = e_s\otimes e_t\) . It then follows from (66) that

\[
\left\{ \begin{array}{l} e _ {s} \llcorner e _ {t} ^ {*} = 0 \quad \text { if   } s \text {   is   not   of   the   form   } t. u \\ e _ {t. u} \llcorner e _ {t} ^ {*} = e _ {u}. \end{array} \right.\tag{72}
\]

Similarly, the symmetric algebra \(\mathbf{S}(\mathbf{M})\) has as basis the set of monomials \(e^{\alpha}\) with \(\alpha \in \mathbf{N}^{n}\) (§ 6, no. 6, Theorem 1) and \(\mathbf{S}(\mathbf{M}^{*})\) the set of monomials \(e^{*\alpha}\) with \(\alpha \in \mathbf{N}^{n}\) ; recall (no. 5) that \(u_{\alpha}\) , for \(|\alpha| = k\) , denotes the element of the basis of \((\mathbf{S}^{k}(\mathbf{M}))^{*}\) , dual to the basis \((e^{\alpha})_{|\alpha| = k}\) of \(\mathbf{S}^{k}(\mathbf{M})\) ; the \(u_{\alpha}\) , for \(\alpha \in \mathbf{N}^{n}\) , therefore form a basis of \(\mathbf{S}(\mathbf{M})^{*\mathrm{gr}}\) . The definition of right inner product by \(e^{\beta}\) in \(\mathbf{S}(\mathbf{M})^{*\mathrm{gr}}\) as the transpose of multiplication by \(e^{\beta}\) in \(\mathbf{S}(\mathbf{M})\) then shows that

\[
\left\{ \begin{array}{l l} u _ {\alpha} \llcorner e ^ {\beta} = 0 & \text { if } \quad \alpha \not \geqslant \beta \\ u _ {\alpha} \llcorner e ^ {\beta} = u _ {\alpha - \beta} & \text { if } \quad \alpha \geqslant \beta . \end{array} \right.\tag{73}
\]

Similarly, since \(\mathbf{S}(\mathbf{M})\) is here canonically identified with the graded dual of \(\mathbf{S}(\mathbf{M})^{*\varepsilon r}\) , \(i(u_{\beta})\) is the graded transpose of multiplication by \(u^{\beta}\) in \(\mathbf{S}(\mathbf{M})^{*\varepsilon r}\) and hence from the multiplication table (33) (no. 5) of the basis \((u_{\alpha})\) we deduce that

\[
\left\{ \begin{array}{l l} e ^ {\alpha} \llcorner u _ {\beta} = 0 & \text { if } \quad \alpha \not \geqslant \beta \\ e ^ {\alpha} \llcorner u _ {\beta} = (\beta , \alpha - \beta) e ^ {\alpha - \beta} & \text { if } \quad \alpha \geqslant \beta . \end{array} \right.\tag{74}
\]

As for the right inner product of an element of \(\mathsf{S}(\mathbf{M})\) by an element of \(\mathsf{S}(\mathbf{M}^*)\) , the definition of this product (no. 9) and formula (34) of no. 5 allow us to deduce from (74) the formulae

\[
\left\{ \begin{array}{l l} e ^ {\alpha} \llcorner e ^ {* \beta} = 0 & \text { if } \quad \alpha \not \geqslant \beta \\ e ^ {\alpha} \llcorner e ^ {* \beta} = \frac {\alpha !}{(\alpha - \beta) !} e ^ {\alpha - \beta} & \text { if } \quad \alpha \geqslant \beta . \end{array} \right.\tag{75}
\]

There are analogous formulae for the inner product of an element of \(\mathsf{S}(\mathbf{M}^{*})\) by an element of \(\mathsf{S}(\mathbf{M})\) interchanging the roles of M and \(M^{*}\) (since \(M^{**}\) is here identified with M).

Remark. Being given the basis \((e_i)_{1\leq i\leq n}\) allows us to identify the algebra \(\mathsf{S}(\mathbf{M})\) with the polynomial algebra \(\mathbf{A}[\mathbf{X}_1,\ldots ,\mathbf{X}_n]\) (§ 6, no. 6); formula (75) shows that the inner product by \(e^{*\alpha}\) is just the differential operator \(\mathbf{D}^{\alpha} = \mathbf{D}_{1}^{\alpha_{1}}\mathbf{D}_{2}^{\alpha_{2}}\ldots \mathbf{D}_{n}^{\alpha_{n}}\) , where \(\mathbf{D}_i = \partial /\partial \mathbf{X}_i\) for \(1\leqslant i\leqslant n\) (§ 10, no. 11, Example).

Consider finally the exterior algebra \(\bigwedge (\mathbf{M})\) , which has as basis the set of elements \(e_{\mathbf{J}}\) , where \(\mathbf{J}\) runs through the set of subsets of the interval \([1, n]\) of \(\mathbf{N}\) (§ 7, no. 8, Theorem 1); similarly \(\bigwedge (\mathbf{M}^*)\) has as basis the elements \(e_{\mathbf{J}}^{*}\) . It follows from formula (68) of no. 9 that

\[
\left\{ \begin{array}{l l} e _ {\mathrm{K}} ^ {*} \lrcorner e _ {\mathrm{J}} = 0 & \text { if } \quad \mathrm{K} \nsubseteq \mathrm{J} \\ e _ {\mathrm{K}} ^ {*} \lrcorner e _ {\mathrm{J}} = (- 1) ^ {p (p - 1) / 2} \rho_ {\mathrm{K}, \mathrm{J} - \mathrm{K}} e _ {\mathrm{J} - \mathrm{K}} & \text { if } \quad \mathrm{K} \subset \mathrm{J} \quad \text { and } \quad p = \operatorname{Card} (\mathrm{K}), \end{array} \right.\tag{76}
\]

where \(\rho_{\mathbf{K},J - \mathbf{K}}\) is the number defined by formula (19) of §7, no. 8. There are analogous formulae with the roles of \(\mathbf{M}\) and \(\mathbf{M}^*\) interchanged.

\subsection{\texorpdfstring{ISOMORPHISMS BETWEEN \(\wedge^{p}(\mathbf{M})\) AND \(\wedge^{n - p}(\mathbf{M}^{*})\) FOR AN \(n\) -DIMENSIONAL FREE MODULE M}{ISOMORPHISMS BETWEEN \textbackslash wedge\^{}\{p\}(\textbackslash mathbf\{M\}) AND \textbackslash wedge\^{}\{n - p\}(\textbackslash mathbf\{M\}\^{}\{*\}) FOR AN n -DIMENSIONAL FREE MODULE M}}

PROPOSITION 12. Let \(\mathbf{M}\) be a free A-module of dimension \(n\) ; let \(e \in \bigwedge^{n}(\mathbf{M})\) be an element forming a basis of \(\bigwedge^{n}(\mathbf{M})\) and let \(e^{*}\) be the element of \(\bigwedge^{n}(\mathbf{M}^{*})\) such that \(\{(-1)^{n(n-1)/2} \theta_{\bigwedge}(e^{*})\}\) is the dual basis of \(\{e\}\) in \((\bigwedge^{n}(\mathbf{M}))^{*}\) . Let \(\phi: \bigwedge(\mathbf{M}^{*}) \to \bigwedge(\mathbf{M})\) be the mapping \(z \mapsto z \lrcorner e^{*}\) and \(\phi': \bigwedge(\mathbf{M}^{*}) \to \bigwedge(\mathbf{M})\) the mapping \(z^{*} \mapsto z^{*} \lrcorner e\) . Let \(\phi_{p}\) (resp. \(\phi_{p}'\) ) be the restriction of \(\phi\) (resp. \(\phi'\) ) to \(\bigwedge^{p}(\mathbf{M})\) (resp. \(\bigwedge^{p}(\mathbf{M}^{*})\) ). Then:

\begin{enumerate}
\def\labelenumi{(\roman{enumi})}
\item
  The mapping \(\phi\) is a left \(\bigwedge (\mathbf{M})\) -module isomorphism and the mapping \(\phi'\) is a left \(\bigwedge (\mathbf{M}^*)\) -module isomorphism; moreover the mappings \(\phi\) and \(\phi'\) are inverses of one another.
\item
  The mapping \(\phi_p\) is an isomorphism of the A-module \(\bigwedge^p (\mathbf{M})\) onto the A-module \(\bigwedge^{n - p}(\mathbf{M}^*)\) and the mapping \(\phi_p'\) is an isomorphism of the A-module \(\bigwedge^p (\mathbf{M}^*)\) onto the A-module \(\bigwedge^{n - p}(\mathbf{M})\) .
\item
  If we write \(\mathbf{B}(u,v^{*}) = \langle u,\theta_{\wedge}(v^{*})\rangle\) for \(u\in \bigwedge (\mathbf{M})\) and \(v^{*}\in \bigwedge (\mathbf{M}^{*})\) then, for \(u^{*}\in \bigwedge^{p}(\mathbf{M}^{*})\) and \(v^{*}\in \bigwedge^{n - p}(\mathbf{M}^{*})\)
\end{enumerate}

\[
\mathbf {B} \left(\phi_ {p} ^ {\prime} (u ^ {*}), v ^ {*}\right) = (- 1) ^ {p (n - p)} \mathbf {B} \left(u ^ {*}, \phi_ {n - p} ^ {\prime} (v ^ {*})\right).\tag{77}
\]

The fact that \(\phi\) is \(\bigwedge (\mathbf{M})\) -linear and \(\phi'\) is \(\bigwedge (\mathbf{M}^*)\) -linear follows from the formulae \((u \wedge v) \lrcorner e^* = u \lrcorner (v \lrcorner e^*)\) and \((u^* \wedge v^*) \lrcorner e = u^* \lrcorner (v^* \lrcorner e)\) (no. 6, formula (37), using the fact that \(\theta_{\wedge}\) is an isomorphism of \(\bigwedge (\mathbf{M}^*)\) onto the opposite algebra to \(\bigwedge (\mathbf{M}^*)\) . On the other hand there exists a basis \((e_i)_{1 \leq i \leq n}\) of \(\mathbf{M}\) such that

\[
e = e _ {1} \wedge e _ {2} \wedge \dots \wedge e _ {n} \quad \text { and } \quad e ^ {*} = (- 1) ^ {n (n - 1) / 2} e _ {1} ^ {*} \wedge e _ {2} ^ {*} \wedge \dots \wedge e _ {n} ^ {*},
\]

where \((e_i^*)\) is the basis dual to \((e_i)\) . We write \(\mathbf{I} = [1, n]\) ; it follows from (76) that, for every subset \(\mathbf{J}\) of \(\mathbf{I}\) with \(p\) elements,

\[
\left\{ \begin{array}{l} \phi (e _ {J}) = (- 1) ^ {n (n - 1) / 2 + p (p - 1) / 2} \rho_ {J, I - J} e _ {I - J} ^ {*} \\ \phi^ {\prime} (e _ {J} ^ {*}) = (- 1) ^ {p (p - 1) / 2} \rho_ {J, I - J} e _ {I - J}. \end{array} \right.\tag{78}
\]

This proves that \(\phi\) and \(\phi'\) are bijective; moreover \(\rho_{J,I-J}\rho_{I-J,J}=(-1)^{p(n-p)}\) ( \(\S 7\) , no. 8, formula (21)); as the number

\[
\frac {n (n - 1)}{2} + \frac {p (p - 1)}{2} + \frac {(n - p) (n - p - 1)}{2} + p (n - p) = n (n - 1)
\]

is even, it follows that \(\phi\) and \(\phi'\) are inverses of one another. Finally, to prove (77), it suffices to take \(u^{*} = e_{J}^{*}\) and \(v^{*} = e_{I - J}^{*}\) ; the verification also follows from the definition of \(\theta_{\wedge}\) , formulae (78) and the relation \(\rho_{J,I - J}\rho_{I - J,J} = (-1)^{p(n - p)}\) (§ 7, no. 8, formula (21)). Note that, for \(u^{*}\in \bigwedge^{p}(\mathbf{M}^{*})\) and \(v^{*}\in \bigwedge^{n - p}(\mathbf{M}^{*})\) , \(\mathbf{B}(\phi_{p}^{\prime}(u^{*}), v^{*})\) is, to within a sign, the coefficient of \(u^{*} \wedge v^{*}\) with respect to the basis \(\{e^{*}\}\) of \(\bigwedge^{n}(\mathbf{M}^{*})\) .

PROPOSITION 13. With the hypotheses and notation of Proposition 11, for every endomorphism g of the A-module M.

\[
(\det g) \phi = \bigwedge (^ {t} g) \circ \phi \circ \bigwedge (g).\tag{79}
\]

Clearly \(\bigwedge(^{t}g) = \theta_{\wedge}^{-1}\circ (^{t}\bigwedge (g))\circ \theta_{\wedge}\) ; since \(\bigwedge (g)\) is an endomorphism of the algebra \(\bigwedge (\mathbf{M})\) and by definition, for all

\[
z \in \bigwedge (\mathbf {M}), \quad \theta_ {\wedge} (\bigwedge (g) (z) \lrcorner e ^ {*}) = \theta_ {\wedge} (e ^ {*}) \llcorner \theta_ {\wedge} (\bigwedge (g) (z)),
\]

we deduce from formula (42) of no. 6 that

\[
\begin{array}{r l} ((\theta_ {\wedge} ^ {- 1} \circ (^ {t} \bigwedge (g)) \circ \theta_ {\wedge}) \circ \phi \circ \bigwedge (g)) (z) & = \theta_ {\wedge} ^ {- 1} (^ {t} \bigwedge (g) (\theta_ {\wedge} (e ^ {*})) \llcorner z) \\ & = z \lrcorner (\bigwedge (t g) (e ^ {*})) = (\det g) (z \lrcorner e ^ {*}) = (\det g) \phi (z) \end{array}
\]

taking account of § 8, no. 4, Proposition 8.

COROLLARY. For every automorphism g of E,

\[
\bigwedge(^ {t} g ^ {- 1}) = (\det g) ^ {- 1} \phi \circ (\bigwedge (g)) \circ \phi^ {- 1}.\tag{80}
\]

\subsection{APPLICATION TO THE SUBSPACE ASSOCIATED WITH A p-VECTOR}

Let K be a field and E a vector space over K. Recall that with every p-vector \(z \in \bigwedge^{p}(\mathbf{E})\) is associated a finite-dimensional subspace \(M_{z}\) of E, namely the smallest vector subspace M of E such that \(z \in \bigwedge^{p}(\mathbf{M})\) (§ 7, no. 2, Corollary to Proposition 4).

PROPOSITION 14. (i) The orthogonal of \(\mathbf{M}_z\) in \(\mathbf{E}^*\) is the set of \(x^{*}\in \mathbf{E}^{*}\) such that \(x^{*}\lrcorner z = 0\) .

\begin{enumerate}
\def\labelenumi{(\roman{enumi})}
\setcounter{enumi}{1}
\tightlist
\item
  The subspace \(\mathbf{M}_z\) associated with \(z\) is the image of \(\bigwedge^{p-1}(\mathbf{E}^*)\) under the mapping \(\lambda_z: u^* \mapsto u^* \lrcorner z\) of \(\bigwedge^{p-1}(\mathbf{E}^*)\) into \(\mathbf{E}\) .
\end{enumerate}

Let \(\mathbf{N}\) denote the image of \(\lambda_z\) . For \(x^* \in \mathbf{E}^*\) and \(u^* \in \bigwedge^{p-1}(\mathbf{E}^*)\) ,

\[
\begin{array}{r l} \langle \theta_ {\wedge} (x ^ {*}), u ^ {*} \lrcorner z \rangle = \langle \theta_ {\wedge} (u ^ {*} \wedge x ^ {*}), z \rangle & = (- 1) ^ {p - 1} \langle \theta_ {\wedge} (x ^ {*} \wedge u ^ {*}), z \rangle \\ & = (- 1) ^ {p - 1} \langle \theta_ {\wedge} (u ^ {*}), x ^ {*} \lrcorner z \rangle . \end{array}
\]

Therefore, for \(x^*\) to be orthogonal to \(\mathbf{N}\) , it is necessary and sufficient that \(x^* \lrcorner z\) be orthogonal to \(\theta_\wedge (\bigwedge (\mathrm{E}^*))\) . Now, the latter condition is equivalent to saying that \(x^* \lrcorner z = 0\) ; for let \((e_\lambda)_{\lambda \in L}\) be a basis of \(\mathbf{E}\) ; giving \(L\) a total ordering, it has been seen (§ 7, no. 8, Theorem 1) that the \(e_j\) , for \(J\) running through the set \(\mathfrak{F}(L)\) of finite subsets of \(L\) , form a basis of \(\bigwedge (\mathrm{E})\) ; it then follows from formula (30) of no. 5 that the elements \(\theta_{\wedge}(e_{\mathrm{J}}^{*})\) are, to within a sign, the coordinate forms on \(\bigwedge(\mathrm{E})\) relative to the basis \((e_{\mathrm{J}})\) ; whence our assertion.

The orthogonal of N therefore consists of the \(x^{*} \in E^{*}\) such that \(x^{*} \lrcorner z = 0\) and the conclusion of (i) will therefore follow from (ii).

We show first that \(\mathbf{N} \subset \mathbf{M}_z\) . Let \(\mathbf{M}\) be a vector subspace of \(\mathbf{E}\) such that \(z \in \bigwedge (\mathbf{M})\) and let \(j: \mathbf{M} \to \mathbf{E}\) be the canonical injection; let \(\mu_z\) denote the mapping \(v^* \mapsto v^* \lrcorner z\) of \(\bigwedge^{p-1}(\mathbf{M}^*)\) into \(\mathbf{M}\) ; it follows from formula (60) of no. 7 that there is a canonical factorization

\[
\lambda_ {z}: \bigwedge^ {p - 1} (\mathrm{E} ^ {*}) \xrightarrow {\bigwedge^ {p - 1} ({} ^ {t} j)} \bigwedge^ {p - 1} (\mathrm{M} ^ {*}) \xrightarrow {\mu_ {z}} \mathrm{M} \xrightarrow {j} \mathrm{E}
\]

which proves that \(\mathbf{N} \subset \mathbf{M}\) and hence \(\mathbf{N} \subset \mathbf{M}_z\) by definition of \(\mathbf{M}_z\) . It remains to verify that \(\mathbf{N} = \mathbf{M}_z\) . Suppose the converse: there would then exist a basis \((e_i)_{1 \leqslant i \leqslant n}\) of \(\mathbf{M}_z\) and an element \(x^* \in \mathbf{E}^*\) such that \(\langle x^*, e_1 \rangle = 1\) , \(\langle x^*, e_j \rangle = 0\) for \(2 \leqslant j \leqslant n\) and such that \(x^*\) is orthogonal to \(\mathbf{N}\) and hence \(x^* \lrcorner z = 0\) . We write \(z = \sum_{\mathbf{H}} a_{\mathbf{H}} e_{\mathbf{H}}\) , where the sum is taken over the subsets of \([1, n]\) with \(p\) elements. By (68) (no. 9),

\[
\begin{array}{l l} x ^ {*} \lrcorner e _ {\mathrm{H}} = 0 & \text { if } \quad 1 \notin \mathrm{H} \\ x ^ {*} \lrcorner e _ {\{1 \} \cup \mathrm{H}} = e _ {\mathrm{H}} & \text { if } \quad \mathrm{H} \subset [ 2, n ] \end{array}
\]

which shows that the relation \(x^{*} \lrcorner z = 0\) implies \(a_{\mathbf{H}} = 0\) for \(1 \in \mathbf{H}\) . But this is impossible, for \(z\) would then belong to \(\bigwedge^{p}(\mathbf{M}')\) , where \(\mathbf{M}'\) is the subspace of \(\mathbf{M}\) generated by \(e_2, \ldots, e_n\) .

\subsection{PURE p-VECTORS. GRASSMANNIANS}

Let K be a field and E a vector space over K. A p-vector \(z \in \bigwedge^{p}(\mathbf{E})\) is called pure (or sometimes decomposable) if it is non-zero and there exist vectors \(x_{1}, \ldots, x_{p}\) in E such that \(z = x_{1} \wedge x_{2} \wedge \cdots \wedge x_{p}\) . For this, it is necessary and sufficient that the subspace \(M_{z}\) associated with z (which is always of dimension \(\geqslant p\) for \(z \neq 0\) ) be exactly of dimension p (since \(\bigwedge^{p}(\mathbf{M}_{z})\) is then of dimension 1). In particular, every non-zero scalar, every non-zero element of \(\mathbf{E} = \bigwedge^{1}(\mathbf{E})\) , every non-zero element of \(\bigwedge^{n}(\mathbf{E})\) , when E is of dimension n, is pure.

PROPOSITION 15. Let E be a vector space of dimension n and let e be an element \(\neq0\) of \(\bigwedge^{n}(\mathbf{E})\) (hence forming a basis of this vector space). Let \(\phi:\bigwedge(\mathbf{E})\to\bigwedge(\mathbf{E}^{*})\) be the vector space isomorphism associated with e (no. 11, Proposition 12). If z is a pure element of \(\bigwedge^{p}(\mathbf{E})\) , then \(\phi(z)\) is a pure element of \(\bigwedge^{n-p}(\mathbf{E}^{*})\) and the subspaces associated with z and \(\phi(z)\) are orthogonal.

The cases \(p = 0\) and \(p = n\) are trivial. Suppose therefore that \(1 \leqslant p \leqslant n - 1\) and let \(z = x_1 \wedge \cdots \wedge x_p \neq 0\) . Then there exists a basis \((e_i)_{1 \leqslant i \leqslant n}\) of \(E\) such that \(e_i = x_i\) for \(1 \leqslant i \leqslant p\) and \(e = e_1 \wedge e_2 \wedge \cdots \wedge e_n\) . It then follows from formula (78) of no. 11 that \(\phi(z) = \pm e_{p+1}^* \wedge \cdots \wedge e_n^*\) , whence the proposition.

COROLLARY. If E is of dimension n, every non-zero \((n - 1)\) -vector E is pure.

PROPOSITION 16. For an element \(z \neq 0\) of \(\bigwedge^{p}(\mathbf{E})\) to be pure, it is necessary and sufficient that, for all \(u^{*} \in \bigwedge^{p-1}(\mathbf{E}^{*})\) ,

\[
(u ^ {*} \lrcorner z) \wedge z = 0.\tag{81}
\]

The case \(p = 0\) is trivial and we assume \(p \geqslant 1\) . If \(z = x_1 \wedge \cdots \wedge x_p\) , formula (68) (no. 9) with \(n = p - 1\) shows that \(u^* \lrcorner z\) is a linear combination of the \(x_i\) ( \(1 \leqslant i \leqslant p\) ), whence (81). If on the other hand the subspace \(M_z\) associated with \(z\) is of dimension \(>p\) , consider a basis \((e_j)_{1 \leqslant j \leqslant n}\) of this subspace with \(n >p\) . It follows from no. 11, Proposition 13 that each of the \(e_j\) is of the form \(u^* \lrcorner z\) for some \(u^* \in \bigwedge^{p-1}(E^*)\) and relation (81) therefore implies \(e_j \wedge z = 0\) for \(1 \leqslant j \leqslant n\) . It follows that in the expression \(z = \sum_{H} a_H e_H\) (where H runs through the set of subsets of {[}1, n{]} with \(p\) elements) all the coefficients \(a_H\) are zero, whence \(z = 0\) , contrary to the hypothesis.

The criterion of Proposition 16 is equivalent to writing conditions (81) when \(u^*\) runs through a basis of \(\bigwedge^{p-1}(\mathbf{E}^*)\) . In particular, suppose that \(\mathbf{E}\) is of finite dimension \(n\) and let \((e_i)_{1\leq i\leq n}\) be a basis of \(\mathbf{E}\) . Conditions (81) are then equivalent to the conditions

\[
(82-(J, H)) \quad \langle e _ {\mathrm{J}} ^ {*}, (e _ {\mathrm{H}} ^ {*} \lrcorner z) \wedge z \rangle = 0
\]

for all subsets J, H of \([1, n]\) such that \(\text{Card}(J) = p + 1\) and \(\text{Card}(H) = p - 1\) . Now, if I and \(I'\) are two subsets of \([1, n]\) with p elements, formulae (76) of no. 10 and multiplication table (20) of § 7, no. 8 show that

\[
\langle e _ {\mathrm{J}} ^ {*}, \langle e _ {\mathrm{H}} ^ {*} \lrcorner e _ {\mathrm{I}}) \wedge e _ {\mathrm{I} ^ {\prime}} \rangle = 0
\]

unless there exists an \(i \in [1, n]\) such that \(\mathbf{I} - \mathbf{H} = \{i\}\) and \(\mathbf{J} - \mathbf{I}' = \{i\}\) , in which case

\[
\langle e _ {\mathrm{J}} ^ {*}, (e _ {\mathrm{H}} ^ {*} \lrcorner e _ {\mathrm{I}}) \wedge e _ {\mathrm{I} ^ {\prime}} \rangle = (- 1) ^ {(p - 1) (p - 2) / 2} \varepsilon_ {i, \mathrm{J,H}}\tag{83}
\]

where \(\varepsilon_{i,J,H} = \rho_{(i),H}\rho_{(i),I'}\) ; it can then be said that for \(i\in J\cap \complement H\) , \(\varepsilon_{i,J,H}\) is equal to \(+1\) if the number of element of \(J\) which are \(< i\) and the number of elements of \(H\) which are \(< i\) have the same parity, and \(-1\) otherwise.

It follows immediately that if we write \(z = \sum_{\mathbf{I}} a_{\mathbf{I}} e_{\mathbf{I}}\) , where \(\mathbf{I}\) runs through the set of subsets of \([1, n]\) with p elements, relation \((82-(J, H))\) is equivalent to the relation

\[
(84 - (\mathrm{J}, \mathrm{H})) \quad \sum_ {i \in J \cup \complement_ {\mathrm{H}}} \varepsilon_ {i, J, \mathrm{H}} a _ {J - \{i \}} a _ {\mathrm{H} \cup \{i \}} = 0.
\]

Relations (84) are called Grassman's relations; these are therefore necessary and sufficient conditions (when J describes the set of subsets with \(p + 1\) elements and H the set of subsets with \(p - 1\) elements of \([1, n]\) ) for an element \(z \neq 0\) of \(\bigwedge^{p}(\mathbf{E})\) to be pure.

Note that relations (84) are not independent. For example, for \(n = 4\) and \(p = 2\) , Grassmann's relations reduce to the single relation

\[
a _ {1 2} a _ {3 4} - a _ {1 3} a _ {2 4} + a _ {1 4} a _ {2 3} = 0.\tag{85}
\]

Let \(D_p(E)\) be the subset of \(\bigwedge^p(E)\) consisting of the pure \(p\) -vectors; clearly \(D_p(E)\) is saturated with respect to the equivalence relation between \(u\) and \(v\) : ``there exists \(\lambda \in K^*\) such that \(v = \lambda u\)'' and two elements \(u, v\) of \(D_p(E)\) are equivalent under this relation if and only if the subspaces \(M_u\) and \(M_v\) of \(E\) which are associated with them are the same. Therefore we thus obtain a canonical bijection of the set of \(p\) -dimensional vector subspaces of \(E\) onto the image \(G_p(E)\) of \(D_p(E)\) in the projective space \(P(\bigwedge^p(E))\) associated with \(\bigwedge^p(E)\) . The subset \(G_p(E)\) of \(P(\bigwedge^p(E))\) is called the Grassmannian of index \(p\) of the vector space \(E\) . When \(E\) is finite-dimensional and \((e_i)_{1 \leq i \leq n}\) is a basis of \(E\) , the Grassmannian of index \(p\) is the set of points of \(P(\bigwedge^p(E))\) for which a system of homogeneous coordinates \((a_I)\) (relative to the basis \((e_I)\) of \(\bigwedge^p(E)\) ) satisfies Grassmann's relations (84).

When \(\mathbf{E} = \mathbf{K}^n\) , we sometimes write \(\mathbf{G}_{n,p}(\mathbf{K})\) instead of \(\mathbf{G}_p(\mathbf{K}^n)\) , so that \(\mathbf{G}_{n,1}(\mathbf{K}) = \mathbf{P}_{n-1}(\mathbf{K})\) . The mapping \(\mathbf{M} \mapsto \mathbf{M}^0\) , which associates with every \(p\) -dimensional subspace of \(\mathbf{K}^n\) the orthogonal subspace in \(\mathbf{E}^*\) (identified with \(\mathbf{K}^n\) under the choice of basis dual to the canonical basis of \(\mathbf{K}^n\) ) therefore defines a canonical bijection of \(\mathbf{G}_{n,p}(\mathbf{K})\) onto \(\mathbf{G}_{n,n-p}(\mathbf{K})\) ; Proposition 15 shows that this bijection is the restriction to \(\mathbf{G}_{n,p}(\mathbf{K})\) of a canonical isomorphism of the projective space \(\mathbf{P}(\bigwedge^p (\mathbf{K}^n))\) onto the projective space \(\mathbf{P}(\bigwedge^{n-p}(\mathbf{K}^n))\) .

\BourbakiAppendixSection{Alternative Algebras. Octonions}

\subsection{ALTERNATIVE ALGEBRAS}

Let A be a commutative ring and F a (not necessarily associative) A-algebra. For any three elements x, y, z of F, we write

\[
a (x, y, z) = x (y z) - (x y) z\tag{1}
\]

(associator of \(x, y, z\) ); \(a\) is obviously an A-trilinear mapping of \(\mathbf{F} \times \mathbf{F} \times \mathbf{F}\) into \(\mathbf{F}\) .

Lemma 1. For all \(p, q, r, s\) in the algebra \(\mathbf{F}\) ,

\[
a (p q, r, s) - a (p, q r, s) + a (p, q, r s) = a (p, q, r) s + p a (q, r, s).\tag{2}
\]

The verification follows immediately from definition (1).

PROPOSITION 1. For an A-algebra F, the following conditions are equivalent:

\begin{enumerate}
\def\labelenumi{(\alph{enumi})}
\item
  For every ordered pair of elements \(x, y\) of \(\mathbf{F}\) , the subalgebra generated by \(x\) and \(y\) is associative.
\item
  The trilinear mapping \((x, y, z) \mapsto a(x, y, z)\) is alternating (§ 7, no. 3).
\item
  For every ordered pair of elements \(x, y\) of \(\mathbf{F}\) , \(x^2y = x(xy)\) and \(yx^2 = (yx)x\) .
\end{enumerate}

Clearly (a) implies (c). We show that (c) implies (b): by definition (§ 7, no. 3) to prove (b), it suffices to verify that \(a(x, x, y) = 0\) and \(a(x, y, y) = 0\) , which is precisely (c).

To prove that (b) implies (a), we use the following 4 lemmas:

Lemma 2. Let E be an A-algebra such that the trilinear mapping \((x,y,z)\mapsto a(x,y,z)\) is alternating, S a generating system of E and U a sub-A-module of E containing S and such that \(sU\subset U\) and \(Us\subset U\) for all \(s\in S\) . Then U = E.

The set \(U'\) of \(x \in E\) such that \(xU \subset U\) and \(Ux \subset U\) is obviously a sub-A-module of E, which contains S by hypothesis. On the other hand, for x, y in \(U'\) and \(u \in U\) , by hypothesis

\[
(x y) u = x (y u) + a (x, y, u) = x (y u) - a (x, u, y) = x (y u) - (x u) y + x (u y) \in \mathrm{U};
\]

on passing to the opposite algebra, we have similarly \(u(xy) \in \mathbf{U}\) . Hence \(\mathbf{U}'\) is a subalgebra of \(\mathbf{E}\) and, since it contains \(\mathbf{S}\) , \(\mathbf{U}' = \mathbf{E}\) . Hence \(\mathbf{EU} \subset \mathbf{U}\) and a fortiori \(\mathbf{UU} \subset \mathbf{U}\) , which proves that \(\mathbf{U}\) is a subalgebra of \(\mathbf{E}\) ; as it contains \(\mathbf{S}\) , \(\mathbf{U} = \mathbf{E}\) , which proves the lemma.

A subset H of F is called strongly associative if \(a(u, v, w) = 0\) when at least two of the elements u, v, w belong to H.

Lemma 3. Suppose that the mapping \(a\) is alternating. If \(\mathbf{H}\) is a strongly associative subset of \(\mathbf{F}\) , the subalgebra of \(\mathbf{F}\) generated by \(\mathbf{H}\) is strongly associative.

As the set of strongly associative subsets of F is inductive, it suffices to prove that if H is a maximal strongly associative subset of F, H is then a subalgebra of F. As H is obviously a sub-A-module of F, it suffices to verify that for any two elements u, v of H, \(H \cup \{uv\}\) is also strongly associative, for by virtue of the definition of H, this will imply \(uv \in H\) . Now, for all \(z \in H\) and all \(t \in F\) , by (2)

\[
a (u v, t, z) - a (u, v t, z) + a (u, v, t z) = 0
\]

since \(\mathbf{H}\) is strongly associative; as \(u, v, z\) are in \(\mathbf{H}\) , also

\[
a (u, v t, z) = a (u, v, t z) = 0,
\]

whence \(a(uv, t, z) = 0\) . Using the fact that a is alternating, this shows that \(a(p, q, r) = 0\) whenever at least two of the elements p, q, r belong to \(H \cup \{uv\}\) whence the lemma.

Lemma 4. Suppose that the mapping \(a\) is alternating. Then, for all \(x \in \mathbf{F}\) , the subalgebra of \(\mathbf{F}\) generated by \(x\) is strongly associative.

\(a(u,v,w) = 0\) whenever two of the three elements \(u, v, w\) are equal to \(x\) and it suffices to apply Lemma 3.

Lemma 5. Suppose that the mapping \(a\) is alternating and let \(\mathbf{X}, \mathbf{Y}\) be two strongly associative subalgebras of \(\mathbf{F}\) . Then the subalgebra of \(\mathbf{E}\) generated by \(\mathbf{X} \cup \mathbf{Y}\) is associative.

Let Z be the set of \(z \in E\) such that \(a(u, v, z) = 0\) for all \(u \in X\) and \(v \in Y\) , this is obviously a sub-A-module containing X and Y since X and Y are strongly associative; by Lemma 2, it will suffice to verify that, for \(u \in X\) and \(v \in Y\) , \(uZ \subset Z\) , \(vZ \subset Z\) , \(Zu \subset Z\) and \(Zv \subset Z\) . Now, for \(u, u'\) in X, \(v \in Y\) and \(z \in Z\) , by (2)

\[
a (u ^ {\prime} u, z, v) - a (u ^ {\prime}, u z, v) + a (u ^ {\prime}, u, z v) = a (u ^ {\prime}, u, z) v + u ^ {\prime} a (u, z, v) = 0
\]

by virtue of the fact that X is strongly associative and the definition of Z. But as X is strongly associative, \(a(u', u, zv) = 0\) and since \(u'u \in X\) , \(a(u'u, z, v) = 0\) by definition of Z. Hence \(a(u', uz, v) = 0\) , which shows that \(uZ \subset Z\) . Applying (2) now with \((p, q, r, s) = (v, z, u, u')\) , we obtain similarly \(Zu \subset Z\) . Interchanging the roles of X and Y and using the fact that a is alternating, we obtain \(vZ \subset Z\) and \(Zv \subset Z\) ; whence the lemma.

It now suffices, to prove that (b) implies (a) in Proposition 1, to take \(\mathbf{X} = \{x\}\) and \(\mathbf{Y} = \{y\}\) , using Lemma 4.

DEFINITION 1. An algebra \(\mathbf{F}\) is called alternative if it satisfies the equivalent conditions of Proposition 1.

An associative algebra is obviously alternative. In no. 3 we shall give an example of an alternative algebra which is not associative.

If \(\mathbf{F}\) is an alternative A-algebra, every A-algebra \(\mathbf{F} \otimes_{\mathbf{A}} \mathbf{A}'\) obtained from \(\mathbf{F}\) by extending the scalars (§ 1, no. 5) is an alternative \(\mathbf{A}'\) -algebra, as follows from condition (b) of Proposition 1.

\subsection{ALTERNATIVE CAYLEY ALGEBRAS}

PROPOSITION 2. Let A be a ring, F a Cayley A-algebra, e its unit element, \(s: x \mapsto \bar{x}\) its conjugation and N: F \(\to\) A its Cayley norm (§ 3, no. 4).

\begin{enumerate}
\def\labelenumi{(\roman{enumi})}
\item
  For \(\mathbf{F}\) to be alternative, it is necessary and sufficient that, for every ordered pair of elements \(x, y\) of \(\mathbf{F}\) , \(x^2y = x(xy)\) .
\item
  If \(\mathbf{F}\) is alternative, then \(\mathbf{N}(xy) = \mathbf{N}(x)\mathbf{N}(y)\) for all \(x,y\) in \(\mathbf{F}\) .
\item
  Suppose that \(\mathbf{F}\) is alternative. For an element \(x \in \mathbf{F}\) to be invertible, it is necessary and sufficient that \(\mathbf{N}(x)\) be invertible in \(\mathbf{Ae}\) ; the inverse of \(x\) is then unique and equal to \(\mathbf{N}(x)^{-1}\bar{x}\) ; denoting this by \(x^{-1}\) ,
\end{enumerate}

\[
x ^ {- 1} (x y) = x \left(x ^ {- 1} y\right) = y
\]

or all \(y \in \mathbf{F}\) .

The condition \(x^{2}y = x(xy)\) is obviously necessary for F to be alternative (no. 1, Proposition 1). Conversely, if it holds for every ordered pair of elements of F, applying it to \(\bar{x}\) and \(\bar{y}\) , it gives \(\bar{x}^{2}\bar{y} = \bar{x}(\bar{x}\bar{y})\) ; applying the conjugation s to this relation, we obtain \(yx^{2} = (yx)x\) , so that conditions (c) of Proposition 1 of no. 1 are satisfied.

Obviously \(a(e, x, y) = 0\) for all \(x, y\) in \(\mathbf{F}\) . If \(\mathbf{F}\) is alternative, we therefore deduce from Proposition 1 (no. 1) that the subalgebra \(\mathbf{G}\) of \(\mathbf{F}\) generated by \(e, x\) and \(y\) is associative. As \(\bar{x} = -x + \mathrm{T}(x) \in -x + \mathrm{Ae}, \bar{x} \in \mathbf{G}\) and similarly \(\bar{y} \in \mathbf{G}\) . Then \(\mathrm{N}(xy) = (xy)(\overline{xy}) = xy.\bar{y}.\bar{x} = \mathrm{N}(y)x\bar{x} = \mathrm{N}(y)\mathrm{N}(x)\) , using the fact that \(\mathrm{N}(y) \in \mathrm{Ae}\) . This proves (ii).

Finally we prove (iii). If \(\mathrm{N}(x)\) is invertible in Ae and we write \(x' = \mathrm{N}(x)^{-1}\bar{x}\) , then \(xx' = x'x = e\) , for \(\mathrm{N}(x) = x\bar{x} = \bar{x}x\) . Conversely if x admits a left inverse \(x''\) , then \(\mathrm{N}(x'')\mathrm{N}(x) = \mathrm{N}(e) = e\) by (ii) and \(\mathrm{N}(x)\) is invertible in Ae; further, as \(x' = \mathrm{N}(x)^{-1}\bar{x}\) is in the subalgebra generated by x and e, the elements \(x, x', x''\) belong to the associative subalgebra generated by x, \(x''\) and e and hence \(x'' = x''(xx') = (x''x)x' = x'\) , whence the uniqueness assertion. The formulae \(x^{-1}(xy) = x(x^{-1}y) = y\) follow from the fact that \(x^{-1}\) , x and y are elements of the subalgebra generated by x, y and e, which is associative.

PROPOSITION 3. Let E be a Cayley A-algebra, \(\gamma\) an element of A and F the Cayley extension of E defined by \(\gamma\) and the conjugation of E (§ 2, no. 5, Proposition 5). For F to be alternative, it is necessary and sufficient that E be associative.

Let \(u = (x, y)\) , \(v = (x', y')\) be two elements of \(\mathbf{F}\) (where \(x, y, x', y'\) are in \(\mathbf{E}\) ). Then (§ 2, no. 5, formula (27))

\[
\left\{ \begin{array}{l} u ^ {2} v = ((x ^ {2} + \gamma \bar {y} y) x ^ {\prime} + \gamma \bar {y} ^ {\prime} (y \bar {x} + y x), (y \bar {x} + y x) \bar {x} ^ {\prime} + y ^ {\prime} (x ^ {2} + \gamma \bar {y} y)) \\ u (u v) = (x (x x ^ {\prime} + \gamma \bar {y} ^ {\prime} y) + \gamma (x ^ {\prime} \bar {y} + \bar {x}. \bar {y} ^ {\prime}) y, y (\bar {x} ^ {\prime} \bar {x} + \gamma \bar {y} y ^ {\prime}) + (y \bar {x} ^ {\prime} + y ^ {\prime} x) x). \end{array} \right.\tag{3}
\]

Using the fact that \(\bar{y}y\) and \(\bar{x} + x\) are in Ae, examining these formulae shows that the associativity of E implies \(u^{2}v = u(uv)\) and hence the fact that F is alternative (Proposition 2). Conversely, if F is alternative, the equation \(u^{2}v = u(uv)\) applied when \(y' = 0\) gives

\[
(y \bar {x} + y x) \bar {x} ^ {\prime} = y (\bar {x} ^ {\prime} \bar {x}) + (y \bar {x} ^ {\prime}) x.
\]

Now the left hand side is equal to \((y\mathrm{T}(x))\bar{x}' = y(\bar{x}'\mathrm{T}(x)) = y(\bar{x}'x + \bar{x}'\bar{x})\) ;

comparing with the right hand side, we obtain \((y\bar{x}^{\prime})x = y(\bar{x}^{\prime}x)\) , which proves the associativity of E, since x, y and \(\bar{x}^{\prime}\) are arbitrary elements of E.

\subsection{OCTONIONS}

Let E be a quaternion algebra of type \((\alpha, \beta, \gamma)\) over A (§2, no. 5, Example 2) and let \(\delta \in A\) . The Cayley extension F of E by \(\delta\) and the conjugation of E is called an octonion algebra over A and is said to be of type \((\alpha, \beta, \gamma, \delta)\) . By Proposition 3 of no. 2, F is an alternative algebra. It has a basis \((e_{i})_{0 \leq i \leq 7}\) of 8 elements, defined by

\[
\begin{array}{l l l l} e _ {0} = (e, 0), & e _ {1} = (i, 0), & e _ {2} = (j, 0), & e _ {3} = (k, 0) \\ e _ {4} = (0, e), & e _ {5} = (0, i), & e _ {6} = (0, j), & e _ {7} = (0, k) \end{array}
\]

where \((e, i, j, k)\) is the basis of E defined loc. cit.; clearly \(e_0\) (also denoted by \(e\) ) is the unit element of F. If \(u = \sum_{i=0}^{7} \xi_i e_i\) is an element of F (with the \(\xi_i \in A\) ), formulae (23), (24) and (31) of § 2, no. 5, give for the conjugate, trace and norm of the octonion \(u\)

\[
\left\{ \begin{array}{c} \bar {u} = (\xi_ {0} + \beta \xi_ {1}) e _ {0} - \sum_ {i = 1} ^ {7} \xi_ {i} e _ {i} \\ \mathrm{T} _ {\mathrm{F}} (u) = 2 \xi_ {0} + \beta \xi_ {1} \\ \mathrm{N} _ {\mathrm{F}} (u) = \xi_ {0} ^ {2} + \beta \xi_ {0} \xi_ {1} - \alpha \xi_ {1} ^ {2} - \gamma (\xi_ {2} ^ {2} + \beta \xi_ {2} \xi_ {3} - \alpha \xi_ {3} ^ {2}) \\ - \delta (\xi_ {4} ^ {2} + \beta \xi_ {4} \xi_ {5} - \alpha \xi_ {5} ^ {2}) + \gamma \delta (\xi_ {6} ^ {2} + \beta \xi_ {6} \xi_ {7} - \alpha \xi_ {7} ^ {2}). \end{array} \right.\tag{4}
\]

Now let \(u = (x, y)\) , \(u' = (x', y')\) and \(u'' = (x'', y'')\) be three octonions (where the elements \(x, x', x'', y, y', y''\) belong to E). Formulae (24) and (27) of § 2, no. 5 give

\[
\begin{array}{r l} & {\mathrm{T} _ {\mathrm{F}} ((u u ^ {\prime}) u ^ {\prime \prime}) = \mathrm{T} (x x ^ {\prime} x ^ {\prime \prime}) + \delta \mathrm{T} (\bar {y} ^ {\prime} y x ^ {\prime \prime}) + \delta \mathrm{T} (\bar {y} ^ {\prime \prime} y \bar {x} ^ {\prime}) + \delta \mathrm{T} (\bar {y} ^ {\prime \prime} y ^ {\prime} x)} \\ & {\mathrm{T} _ {\mathrm{F}} (u (u ^ {\prime} u ^ {\prime \prime})) = \mathrm{T} (x x ^ {\prime} x ^ {\prime \prime}) + \delta \mathrm{T} (x ^ {\prime \prime} \bar {y} ^ {\prime} y) + \delta \mathrm{T} (\bar {x} ^ {\prime} \bar {y} ^ {\prime \prime} y) + \delta \mathrm{T} (x \bar {y} ^ {\prime \prime} y ^ {\prime})} \end{array}
\]

(where T denotes trace in E and use is made of the fact that E is associative). As \(\mathrm{T}(xy) = \mathrm{T}(yx)\) for all quaternions x, y (§ 2, no. 4, formula (17)), it follows that

\[
\mathrm{T} _ {\mathrm{F}} ((u u ^ {\prime}) u ^ {\prime \prime}) = \mathrm{T} _ {\mathrm{F}} (u (u ^ {\prime} u ^ {\prime \prime})).\tag{5}
\]

We study in particular octonions of type \((-1, 0, -1, -1)\) ; formulae (4) then simplify to

\[
\left\{ \begin{array}{c} {\bar {u}} = \xi_ {0} e _ {0} - \sum_ {i = 1} ^ {7} \xi_ {i} e _ {i} \\ \mathrm{T} _ {\mathrm{F}} (u) = 2 \xi_ {0} \\ \mathrm{N} _ {\mathrm{F}} (u) = \sum_ {i = 0} ^ {7} \xi_ {i} ^ {2}. \end{array} \right.\tag{6}
\]

*If we take A to be the field R of real numbers, the octonions of type \((-1,0,-1,-1)\) over R are called Cayley octonions (or octaves). It follows from Proposition 2 (ii) of no. 2 that every Cayley octonion \(\neq0\) is invertible.*

PROPOSITION 4. Let F be an octonion algebra of type \((-1,0,-1,-1)\) over A. There exists a vector space V of dimension 3 over the field with two elements Z/2Z and a bijection \(\lambda\mapsto e_{\lambda}^{\prime}\) of V onto the basis \((e_{i})_{0\leq i\leq7}\) such that

\[
e _ {0} ^ {\prime} = e _ {0}, \quad e _ {\lambda} ^ {\prime} e _ {\mu} ^ {\prime} = \pm e _ {\lambda + \mu} ^ {\prime}\tag{7}
\]

for all \(\lambda, \mu\) in V. In order that \(e_{\lambda}'(e_{\mu}'e_{\nu}') = (e_{\lambda}'e_{\mu}')e_{\nu}'\) , it suffices that, in V, \(\lambda, \mu, \nu\) be linearly dependent over Z/2Z; this condition is necessary if \(2 \neq 0\) in A.

We preserve the notation at the beginning of this no. It follows from formulae (33) of §2, no. 5 that the set S consisting of the elements \(\pm e_0\) , \(\pm e_1\) , \(\pm e_2\) , \(\pm e_3\) is stable under multiplication. Moreover, for \(x, y, y'\) in E, by formula (22) of §2, no. 5,

\[
(x, 0) (0, y ^ {\prime}) = (0, y ^ {\prime} x), \quad (0, y ^ {\prime}) (x, 0) = (0, y ^ {\prime} \bar {x}), \quad (0, y) (0, y ^ {\prime}) = (- \bar {y} ^ {\prime} y, 0)\tag{8}
\]

so that the set T consisting of the elements \(\pm e_{i}\) ( \(0 \leqslant i \leqslant 7\) ) is stable under multiplication; moreover, its multiplication table is independent of the ring A.

In particular, let \(\mathbf{A}''\) be the field \(\mathbf{Z}/2\mathbf{Z}\) with two elements and let \(\mathbf{E}''\) be the quaternion algebra of type (1, 0, 1) over \(\mathbf{A}''\) and \(\mathbf{F}''\) the algebra of octonions of type (1, 0, 1, 1) over \(\mathbf{A}''\) ; let \((e_i'')_{0 \leq i \leq 7}\) be the basis of \(\mathbf{F}''\) formed as described above. Since \(-e_i'' = e_i''\) , the set \(\mathbf{T}''\) of \(e_i''\) has 8 elements and is stable under multiplication; moreover, it follows immediately from the above that the mapping \(\theta: \mathbf{T} \to \mathbf{T}''\) such that \(\theta(e_i) = \theta(-e_i) = e_i''\) for \(0 \leq i \leq 7\) is a homomorphism for multiplication. Moreover the quaternion algebra \(\mathbf{E}''\) is in this case commutative and hence \(\mathbf{F}''\) is associative ( \(\S 2\) , no. 5, Proposition 5); further, conjugation in \(\mathbf{F}''\) is in this case the identity. Hence \(\mathbf{T}''\) is a group and formulae (8) show that it is commutative; these formulae and formulae (33) of \(\S 2\) , no. 5 show that the square of every element of \(\mathbf{T}''\) is the unit. If V is used to denote the group \(\mathbf{T}''\) written additively, V can be given a unique vector space structure over \(\mathbf{Z}/2\mathbf{Z}\) , necessarily of dimension 3 since

\[
\operatorname{Card} (\mathbf {V}) = 8 = 2 ^ {\dim (\mathbf {V})}.
\]

For all \(\lambda \in V\) , let \(e_{\lambda}'\) then denote the element of \((e_i)_{0 \leq i \leq 7}\) such that \(\theta(e_{\lambda}') = \dot{\lambda}\) ; then \(e_0' = e_0\) ; moreover, as \(\theta\) is a homomorphism and the relation \(\theta(x) = \theta(y)\) is equivalent to \(x = \pm y\) , \(e_{\lambda}' e_{\mu}' = \pm e_{\lambda + \mu}'\) . If \(\lambda, \mu, \nu\) are linearly independent over \(\mathbf{Z}/2\mathbf{Z}\) , they form a basis of \(V\) and hence all the elements \(e_i\) ( \(0 \leqslant i \leqslant 7\) ) would belong to the subalgebra generated by \(e_{\lambda}', e_{\mu}'\) and \(e_{\nu}'\) ; when \(2 \neq 0\) in \(A\) , \(e_{\lambda}'(e_{\mu}' e_{\nu}') = (e_{\lambda}' e_{\mu}') e_{\nu}'\) is therefore impossible, for \(F\) would be associative and hence \(E\) would be commutative ( \(\S 2\) , no. 5, Proposition 5), which contradicts relations (33) of \(\S 2\) , no. 5. On the other hand, if \(\lambda, \mu, \nu\) are linearly

dependent in V, the three elements \(e_{\lambda}^{\prime}\) , \(e_{\mu}^{\prime}\) , \(e_{\nu}^{\prime}\) belong to a subalgebra with 2 generators of F, which is therefore associative (no. 1, Proposition 1); whence the conclusion.

Remark. As \(\bar{e}_{\lambda}^{\prime} = -e_{\lambda}^{\prime}\) for \(\lambda \neq 0\) ,

\[
e _ {\lambda} ^ {\prime 2} = - e _ {0} \quad \text { for } \lambda \neq 0,
\]

\[
e _ {\mu} ^ {\prime} e _ {\lambda} ^ {\prime} = - e _ {\lambda} ^ {\prime} e _ {\mu} ^ {\prime} \quad \text { for } \lambda \neq 0, \mu \neq 0 \text { and } \mu \neq \lambda .
\]

\BourbakiExercisesHeading

\BourbakiExercisesSection

\begin{enumerate}
\def\labelenumi{\arabic{enumi}.}
\tightlist
\item
  Let E be an algebra over a (commutative) field K and L a commutative extension field of K. Show that, if the algebra \(\mathrm{E}_{(L)}\) admits a unit element, so does E (cf.~II, §3, no. 5, Proposition 6).
\end{enumerate}

\BourbakiExercisesSection

\begin{enumerate}
\def\labelenumi{\arabic{enumi}.}
\tightlist
\item
  Let E be an A-algebra with a basis of two elements \(e_1, e_2\) and a unit element \(e = \lambda e_1 + \mu e_2\) .
\end{enumerate}

\begin{enumerate}
\def\labelenumi{(\alph{enumi})}
\item
  Show that the ideal \(a\) of A generated by \(\lambda\) and \(\mu\) is the whole ring A (in the contrary case, note that \(\mathrm{E} / \mathrm{aE} = \mathrm{E} \otimes_{\mathbf{A}} (\mathrm{A} / \mathrm{a}) = \{0\}\) would hold contrary to the fact that \(\mathrm{A} / \mathrm{a} \neq \{0\}\) and E is a free A-module).
\item
  If \(\alpha, \beta\) are two elements of A such that \(\alpha\lambda + \beta\mu = 1\) , show that the elements \(e\) and \(u = \beta e_1 - \alpha e_2\) also form a basis of E and deduce that E is a quadratic A-algebra.
\end{enumerate}

\begin{enumerate}
\def\labelenumi{\arabic{enumi}.}
\setcounter{enumi}{1}
\item
  Show that every quadratic algebra over \(\mathbf{Z}\) is isomorphic to an algebra of type \((0, n)\) with \(n \in \mathbf{N}\) , or an algebra of type \((m, 0)\) with \(m\) a non-square, or an algebra of type \((m, 1)\) , where \(m\) is an integer which is not of the form \(k(k - 1)\) ; moreover, no two of these algebras are isomorphic.
\item
  \begin{enumerate}
  \def\labelenumii{(\alph{enumii})}
  \tightlist
  \item
    Let K be a commutative field. Show that, for K to be the centre of a quaternion algebra of type \((\alpha, \beta, \gamma)\) over K, it is necessary and sufficient that one of the following cases hold: (1) \(\beta\gamma \neq 0\) ; (2) \(\gamma = 0\) , \(\beta \neq 0\) and \(\beta^{2} + 4\alpha \neq 0\) ; (3) \(\beta = 0\) , \(\alpha \neq 0\) or \(\gamma \neq 0\) and \(2 \neq 0\) in K.
  \end{enumerate}
\end{enumerate}

\begin{enumerate}
\def\labelenumi{(\alph{enumi})}
\setcounter{enumi}{1}
\item
  Let K be a commutative field such that \(2 \neq 0\) in K; show that the quaternion algebra over K of type \((1, \gamma)\) is isomorphic to the matrix algebra \(\mathbf{M}_{2}(\mathbf{K})\) (consider the basis of this algebra consisting of the elements \(\frac{1}{2}(1 + i)\) , \(\frac{1}{2}(1 - i)\) , \((j + k/2\gamma)\) , \(\frac{1}{2}(j - k)\) ).
\item
  Let K be a commutative field such that \(2 \neq 0\) in K and let E be the quaternion algebra over K of type \((\alpha, \gamma)\) . A quaternion \(z \in E\) is called pure if \(\overline{z} = -z\) (or also \(T(z) = 0\) ). Show that if z is pure, so is every quaternion \(tzt^{-1}\) , where \(t \in E\) is invertible.
\item
  If \(u, v\) are any two quaternions, \(uv - vu\) is pure. Deduce that if \(u, v, w\) are any three quaternions, then
\end{enumerate}

\[
(u v - v u) ^ {2} w = w (u v - v u) ^ {2}.
\]

\begin{enumerate}
\def\labelenumi{(\alph{enumi})}
\setcounter{enumi}{4}
\tightlist
\item
  Show that if there exists in E a quaternion \(z \neq 0\) such that \(\mathrm{N}(z) = 0\) , there also exists a pure quaternion \(z' \neq 0\) such that \(\mathrm{N}(z') = 0\) . (Note, in the notation of formula (33) (no. 5), that, if \(\alpha\) is not a square in K, the existence of \(z \neq 0\) in E such that \(\mathrm{N}(z) = 0\) is equivalent to the existence of an element \(y \in \mathrm{K} + \mathrm{Ki}\) such that \(\gamma = \mathrm{N}(y)\) .)
\end{enumerate}

\begin{enumerate}
\def\labelenumi{\arabic{enumi}.}
\setcounter{enumi}{4}
\tightlist
\item
  Over a commutative field K in which 2 = 0, every quaternion algebra E of type \((\alpha, 0, \gamma)\) is commutative and the square of every \(x \in E\) belongs to K; the subalgebra \(\mathrm{K}(x)\) of E generated by an element \(x \in K\) is therefore a quadratic algebra of type \((\lambda, 0)\) over K and E is a quadratic algebra over \(\mathrm{K}(x)\) .
\end{enumerate}

Show that the set C of \(x \in E\) whose square is equal to the square of an element of K is a vector subspace of E whose dimension is equal to 1, 2 or 4. If C is of dimension 1 (in which case C = K), E is a field. If C is of dimension 2, there exists a quadratic algebra \(\mathrm{K}(x)\) contained in E which is a field and E has a basis over \(\mathrm{K}(x)\) consisting of two elements 1 and u with \(u^{2} = 0\) ; the set of \(y \in E\) such that \(y^{2} = 0\) is an ideal a of dimension 2 over K and E/a is isomorphic to \(\mathrm{K}(x)\) . Finally, if C is of dimension 4 (and hence equal to E), the set a of \(y \in E\) such that \(y^{2} = 0\) is an ideal of dimension 3 and E/a is isomorphic to K; there exists in E a basis \((1, e_{1}, e_{2}, e_{3})\) such that \(e_{1}^{2} = e_{2}^{2} = e_{3}^{2} = 0\) , \(e_{1}e_{2} = e_{3}\) , \(e_{1}e_{3} = e_{2}e_{3} = 0\) ; \(Ke_{3} = b\) is the only ideal of dimension 1 in E, b is the annihilator of a and a/b is the direct sum of two ideals of E/b, of dimension 1, which annihilate one another.

\begin{enumerate}
\def\labelenumi{\arabic{enumi}.}
\setcounter{enumi}{5}
\tightlist
\item
  Let \(\mathbf{K}\) be a commutative field in which \(2 \neq 0\) and \(\mathbf{E}\) a quaternion algebra over \(\mathbf{K}\) of type \((0, 0, \gamma)\) .
\end{enumerate}

\begin{enumerate}
\def\labelenumi{(\alph{enumi})}
\item
  If \(\gamma\) is not a square in \(\mathbf{K}\) , there exists no left (resp. right) ideal in \(\mathbf{E}\) of dimension 1 over \(\mathbf{K}\) ; the set \(a\) of \(x \in \mathbf{E}\) such that \(x^2 = 0\) is a two-sided ideal of dimension 2 over \(\mathbf{K}\) and \(\mathbf{E} / a\) is a field, a quadratic algebra over \(\mathbf{K}\) .
\item
  If \(\gamma\) is a square \(\neq 0\) in \(\mathbf{K}\) , there exists in \(\mathbf{E}\) a basis \((e_i)_{1\leq i\leq 4}\) with the following multiplication table:
\end{enumerate}

\(e_1\)

\(e_2\)

\(e_3\)

\(e_4\)

\(e_1\)

\(e_1\)

0

\(e_3\)

0

\(e_2\)

0

\(e_2\)

0

\(e_4\)

\(e_3\)

0

\(e_3\)

0

0

\(e_4\)

\(e_4\)

0

0

0

The set a of \(x \in E\) such that \(x^{2} = 0\) is a two-sided ideal of dimension 2, the direct sum of the two two-sided ideals \(Ke_{3}\) and \(Ke_{4}\) which annihilate one another; the latter are the only (left or right) ideals of dimension 1 in E. The quotient algebra E/a is isomorphic to the product of two fields isomorphic to K.

\begin{enumerate}
\def\labelenumi{(\alph{enumi})}
\setcounter{enumi}{2}
\tightlist
\item
  If \(\gamma = 0\) , the set \(a\) of \(x \in E\) such that \(x^2 = 0\) is a two-sided ideal of dimension 3; \(Kk = b\) is the only (left or right) ideal of dimension 1 in \(E\) ; it is a two-sided ideal, the left and right annihilator of \(a\) . In the quotient algebra \(E/b\) , \(a/b\) is the direct sum of two two-sided ideals of dimension 1 which annihilate one another; finally \(E/a\) is a field isomorphic to \(K\) .
\end{enumerate}

\begin{enumerate}
\def\labelenumi{\arabic{enumi}.}
\setcounter{enumi}{6}
\tightlist
\item
  Let K be a commutative field in which \(2 \neq 0\) and E the algebra over K with a basis of four elements 1, i, j, k (1 the unit element) with multiplication table
\end{enumerate}

\[
i ^ {2} = j ^ {2} = k ^ {2} = 1, \quad i j = j i = k, \quad j k = k j = i, \quad k i = i k = j.
\]

Show that E is isomorphic to the product of four fields isomorphic to K (consider the basis of E consisting of the elements \((1 + \varepsilon i)(1 + \varepsilon'j)\) , where \(\varepsilon\) and \(\varepsilon'\) are equal to 1 or -1).

The algebra E is the algebra (over K) of the product of two cyclic groups of order 2. Generalize this to the algebra of the product group of n cyclic groups of order 2 (cf.~§ 4, no. 1, Example 2).

\begin{enumerate}
\def\labelenumi{\arabic{enumi}.}
\setcounter{enumi}{7}
\item
  The quaternionic group \(\mathfrak{Q}\) (I, § 6, Exercise 4) is isomorphic to the multiplicative group of the eight quaternions \(\pm 1\) , \(\pm i\) , \(\pm j\) , \(\pm k\) in the quaternion algebra of type \((-1, -1)\) over a field in which \(2 \neq 0\) . Show that the algebra of the group \(\mathfrak{Q}\) over a field K in which \(2 \neq 0\) is isomorphic to the product of four fields isomorphic to K and the quaternion algebra of type \((+1, -1)\) over K (if \(c\) is the element of \(\mathfrak{Q}\) corresponding to the quaternion \(-1\) , the elements of \(\mathfrak{Q}\) can be written as \(e\) , \(i, j, k, c, ci, cj, ck\) ; consider the basis of E consisting of the elements \(\frac{1}{2}(e + c)\) , \(\frac{1}{2}(e - c)\) , \(\frac{1}{2}(e + c)i\) , \(\frac{1}{2}(e - c)i\) , \(\frac{1}{2}(e + c)j\) , \(\frac{1}{2}(e - c)j\) , \(\frac{1}{2}(e + c)k\) , \(\frac{1}{2}(e - c)k\) ).
\item
  \begin{enumerate}
  \def\labelenumii{(\alph{enumii})}
  \tightlist
  \item
    Show that the algebra E of the dihedral group \(\mathbf{D}_4\) of order 8 (I, §6, Exercise 4) over a field K in which \(2 \neq 0\) is isomorphic to the product of four fields isomorphic to K and the algebra of matrices of order 2 over K. (If \(a, b\) are the two generators of \(\mathbf{D}_4\) introduced in I, §6, Exercise 4, the elements of \(\mathbf{D}_4\) can be written as \(a^i b^j\) with \(0 \leqslant i \leqslant 3\) , \(0 \leqslant j \leqslant 1\) ; consider the basis of E consisting of the four elements \(\frac{1}{2}(e + a^2)\) , \(\frac{1}{2}(e - a^2)\) , \(\frac{1}{2}(a + a^3)\) , \(\frac{1}{2}(a - a^3)\) and these four elements multiplied on the right by \(b\) ; use Exercise 4.) Deduce that, if \(-1\) is a square in K, the algebras over K of the non-isomorphic groups \(\mathfrak{Q}\) and \(\mathbf{D}_4\) are isomorphic.
  \end{enumerate}
\end{enumerate}

\begin{enumerate}
\def\labelenumi{(\alph{enumi})}
\setcounter{enumi}{1}
\tightlist
\item
  Show similarly that the algebra F of the dihedral group \(D_{3}\) of order 6 over a field K in which \(2 \neq 0\) and \(3 \neq 0\) is isomorphic to the product of two fields isomorphic to K and the matrix algebra of order 2 over K (consider here the basis of F consisting of the elements \(e + a + a^{2}\) , \(a + a^{2} - 2e\) , \(a - a^{2}\) and these three elements multiplied on the right by b).
\end{enumerate}

\begin{enumerate}
\def\labelenumi{\arabic{enumi}.}
\setcounter{enumi}{9}
\item
  Let G be a group and H a normal subgroup of G. Show that, if a is the two-sided ideal of the group algebra \(\mathbf{A}^{(\mathbf{G})}\) , generated by the elements ts - s, where t runs through H and s runs through G, the group algebra \(\mathbf{A}^{(\mathbf{G}/\mathbf{H})}\) is isomorphic to the quotient algebra \(\mathbf{A}^{(\mathbf{G})}/\mathbf{a}\) .
\item
  Let A be a commutative ring and S a commutative monoid with identity element e. Suppose that an external law \((s, x) \mapsto x^{s}\) is given on A with S as domain of operators, such that, for all \(s \in S\) , the mapping \(x \mapsto x^{s}\) is an automorphism of the ring A and \((x^{s})^{t} = x^{st}\) for all \(x \in A\) and s and t in S (which implies that e is the identity operator for the external law considered). Then a multiplicative internal law of composition is defined on the A-module \(\mathrm{A}^{(\mathrm{S})}\) , whose canonical basis is denoted by \((b_{\mathrm{s}})\) , by the relation
\end{enumerate}

\[
\left(\sum_ {s} x _ {s} b _ {s}\right) \left(\sum_ {s} y _ {s} b _ {s}\right) = \sum_ {s} \left(\sum_ {t u = s} a _ {t, u} x _ {t} y _ {u} ^ {s}\right) b _ {s}
\]

where the \(a_{s,t}\) belong to A.

Show that this law and addition on \(\mathbf{A}^{(\mathbf{S})}\) define a ring structure on this set, provided the \(a_{s,t}\) satisfy the conditions

\[
a _ {s, t} a _ {s t, u} = a _ {t u} ^ {s} a _ {s, t u}
\]

for all s, t, u in S. The element \(b_{e}\) is unit element of this ring if also

\[
a _ {e, s} = a _ {s, e} = 1
\]

for all \(s \in S\) ; in that case A can be identified with a subring of \(A^{(S)}\) and if C denotes the subring of A consisting of the elements invariant under all the automorphisms \(x \mapsto x^s\) , \(A^{(S)}\) is a C-algebra, called the crossed product of the ring A and the monoid S, relative to the factor system \((a_{s,t})\) . If the factor system \((a_{s,t})\) is replaced by \(\left(\frac{c_s c_t^s}{c_{st}} a_{s,t}\right)\) , where, for all \(s \in S\) , \(c_s\) is an invertible element of A and \(c_e = 1\) , the new crossed product obtained is isomorphic to the crossed product defined by the factor system \((a_{s,t})\) .

If in particular A is taken to be a quadratic algebra over a ring C and S a cyclic group of order 2, say \(\{e, s\}\) , such that \(x^s = \bar{x}\) for \(x \in \mathbf{A}\) , show that every crossed product of A and S is a quaternion algebra over C.

\begin{enumerate}
\def\labelenumi{\arabic{enumi}.}
\setcounter{enumi}{11}
\tightlist
\item
  Let I be a non-empty set. Show that the elements of the free magma M(I) can be identified with the sequences \((a_{i})_{1 \leqslant i \leqslant n}\) (n an arbitrary integer \(\geqslant 1\) ) where the \(a_{i}\) are either elements of I or a special symbol P (representing the ``open brackets''), where these sequences satisfy the following conditions:
\end{enumerate}

\begin{enumerate}
\def\labelenumi{(\roman{enumi})}
\item
  The length \(n\) of the sequence is equal to twice the number of symbols \(\mathbf{P}\) which appear in it, plus 1.
\item
  For all \(k \leqslant n - 1\) , the number of symbols \(\mathbf{P}\) which appear in the subsequence \((a_i)_{1 \leqslant i \leqslant k}\) is \(\geqslant k / 2\) .
\end{enumerate}

(Use Set Theory, I, Appendix to associate with each sequence \((a_i)\) an element of M(I); for example, to the sequence PPPxyzPxy corresponds the word (((xy)z)(xy)).)

¶ 13. Let A be a commutative ring and E an A-module. An A-module \(E_{n}\) is defined inductively by writing

\[
\mathbf {E} _ {1} = \mathbf {E}, \quad \mathbf {E} _ {n} = \bigoplus_ {p + q = n} \left(\mathbf {E} _ {p} \otimes_ {\mathbf {A}} \mathbf {E} _ {q}\right) \quad \text { for } n \geqslant 2.
\]

We write \(\mathbf{ME} = \bigoplus_{n=1}^{\infty} \mathbf{E}_n\) and an A-algebra structure is defined on ME by setting, for \(x_p \in \mathbf{E}_p\) and \(x_q \in \mathbf{E}_q\) , \(x_p x_q = x_p \otimes x_q \in \mathbf{E}_{p+q}\) .

\begin{enumerate}
\def\labelenumi{(\alph{enumi})}
\item
  Show that for every A-algebra B and every A-linear mapping \(f: \mathbf{E} \to \mathbf{B}\) , there exists one and only one A-homomorphism of algebras \(\mathbf{ME} \to \mathbf{B}\) extending \(f\) . ME is called the free algebra of the A-module E.
\item
  An integer \(a(n)\) is defined by induction on \(n\) by the formulae
\end{enumerate}

\[
a (1), \quad a (n) = \sum_ {p + q = n} a (p) a (q) \quad \text { if } n \geqslant 2.
\]

Show that \(E_{n}\) is isomorphic to the direct sum of \(a(n)\) modules isomorphic to \(E^{\otimes n}\) .

\begin{enumerate}
\def\labelenumi{(\alph{enumi})}
\setcounter{enumi}{2}
\tightlist
\item
  Let \(f(\mathbf{T})\) be the formal power series \(\sum_{n=1}^{\infty} a(n)\mathbf{T}^n\) . Show that
\end{enumerate}

\[
f (\mathbf {T}) = \mathbf {T} + (f (\mathbf {T})) ^ {2}.
\]

*Deduce the formulae

\[
f (\mathrm{T}) = \frac {1 - \sqrt {1 - 4 \mathrm{T}}}{2} \quad \text { and } \quad a (n) = 2 ^ {n - 1} \frac {1 . 3 . 5 \dots (2 n - 1)}{n !}. *
\]

\begin{enumerate}
\def\labelenumi{(\alph{enumi})}
\setcounter{enumi}{3}
\tightlist
\item
  If I is any set and \(E = A^{(I)}\) , show that \(\operatorname{Lib}_{\mathbf{A}}(\mathbf{I})\) is identified with ME. There are canonical homomorphisms \(\mathbf{M}(\mathbf{I}) \xrightarrow{w} \mathbf{N}^{(\mathbf{I})} \xrightarrow{l} \mathbf{N}\) whose composition is length in \(\mathbf{M}(\mathbf{I})\) ; \(\omega(m)\) is called the multidegree of \(m \in \mathbf{M}(\mathbf{I})\) ; the mapping \(\omega\) defines on \(\operatorname{Lib}_{\mathbf{A}}(\mathbf{I})\) a graduation of type \(N^{(I)}\) . Show that the set \((\operatorname{Lib}_{\mathbf{A}}(\mathbf{I}))_{\alpha}\) of homogeneous elements of multidegree \(\alpha \in N^{(I)}\) under this graduation is a free A-module of rank equal to
\end{enumerate}

\[
a (\alpha) = a (n) \frac {n !}{\alpha !} = 2 ^ {n - 1} \frac {1 . 3 . 5 \dots (2 n - 1)}{\alpha !}
\]

where \(\alpha! = \prod_{i \in I} (\alpha(i))!\) and \(n = \sum_{i \in I} \alpha(i) = l(\alpha)\) is the length of \(\alpha\) .

\begin{enumerate}
\def\labelenumi{\arabic{enumi}.}
\setcounter{enumi}{13}
\tightlist
\item
  \begin{enumerate}
  \def\labelenumii{(\alph{enumii})}
  \tightlist
  \item
    Let M be a monoid whose law of composition is denoted by \(\top\) ; suppose further that M is totally ordered by an order relation \(x \leqslant y\) such that the relations \(x < y\) and \(x' \leqslant y'\) , or \(x \leqslant y\) and \(x' < y'\) , imply \(x \top x' < y \top y'\) . Show that if A is an integral domain, the algebra over A of the monoid M has no divisors of 0.
  \end{enumerate}
\end{enumerate}

\begin{enumerate}
\def\labelenumi{(\alph{enumi})}
\setcounter{enumi}{1}
\item
  Deduce that if \(\mathbf{A}\) is an integral domain, every polynomial algebra \(\mathbf{A}[(\mathbf{X}_i)_{i\in I}]\) is an integral domain.
\item
  If A is an integral domain and f and g two polynomials in \(\mathrm{A}[(\mathrm{X}_{i})_{i\in\mathbf{I}}]\) such that fg is homogeneous and \(\neq0\) , show that f and g are homogeneous. In particular, the invertible elements of the ring \(\mathrm{A}[(\mathrm{X}_{i})_{i\in\mathbf{I}}]\) are the invertible elements of the ring A.
\end{enumerate}

\begin{enumerate}
\def\labelenumi{\arabic{enumi}.}
\setcounter{enumi}{14}
\item
  Let A be a commutative ring and n an integer \(\geqslant2\) such that for all \(a\in A\) the equation nx=a has a solution \(x\in A\) . Let m be an integer \(\geqslant1\) and u a polynomial in A{[}X{]} of the form \(X^{mn}+f(X)\) where f is of degree \(\leqslant mn-1\) ; show that there exists a polynomial \(v\in A[X]\) of the form \(X^{m}+w\) , where w is of degree \(\leqslant m-1\) , such that \(u-v^{n}\) is zero or of degree \(<m(n-1)\) .
\item
  Let A be a commutative ring and \(u = \sum_{k=0}^{m} a_k X^k\) a divisor of 0 in the ring A{[}X{]}. Show that if \(v = \sum_{k=0}^{n} b_k X^k\) is an element \(\neq 0\) of A{[}X{]} of degree \(n\) such that \(uv = 0\) , there exists a polynomial \(w \neq 0\) of degree \(\leqslant n - 1\) such that \(uw = 0\) . (Reduce it to the case where \(b_0 \neq 0\) ; if \(a_k v = 0\) for \(0 \leqslant k \leqslant m - 1\) , show that we can take \(w = b_0\) ; if \(a_k v = 0\) for
\end{enumerate}

\[
0 \leqslant k <   p \leqslant m - 1
\]

and \(a_{p}v \neq 0\) , show that \(a_{p}b_{0} = 0\) and therefore we can take \(w = \sum_{k=0}^{n-1} a_{p}b_{k+1}\mathbf{X}^{k}\) . Deduce that there exists an element \(c \neq 0\) of \(\mathbf{A}\) such that \(cu = 0\) .

\begin{enumerate}
\def\labelenumi{\arabic{enumi}.}
\setcounter{enumi}{16}
\item
  Let A be a commutative ring. Show that, in the algebra A{[}X{]} of polynomials in one indeterminate over A, the mapping \((u,v)\mapsto u(v)\) is a law of composition which is associative and left distributive with respect to addition and multiplication in A{[}X{]}. If A is an integral domain, show that the relation \(u(v)=0\) implies that u=0 or that v is constant and that, if \(u\neq0\) and the degree of v is \textgreater0, the degree of \(u(v)\) is equal to the product of the degrees of u and v.
\item
  Let K be a commutative field and K{[}X{]} the ring of polynomials in one indeterminate over K. Show that every automorphism s of the ring K{[}X{]} leaves K invariant (Exercise 14(c)) and induces on K an automorphism of this field; show further that necessarily \(s(\mathbf{X}) = a\mathbf{X} + b\) , with \(a \neq 0\) in K and \(b \in \mathbf{K}\) (cf.~Exercise 17). Conversely, show that being given an automorphism \(\sigma\) of \(\mathbf{K}\) and two elements \(a, b\) of \(\mathbf{K}\) such that \(a \neq 0\) defines one and only one automorphism \(s\) of \(\mathbf{K}[\mathbf{X}]\) such that \(s|_{\mathbf{K}} = \sigma\) and \(s(\mathbf{X}) = a\mathbf{X} + b\) . If \(G\) is the group of all automorphisms of the ring \(\mathbf{K}[\mathbf{X}]\) and \(N\) the subgroup of \(G\) consisting of the K-algebra structure of \(\mathbf{K}[\mathbf{X}]\) , show that \(N\) is a normal subgroup of \(G\) and that \(G / N\) is isomorphic to the automorphism group of the field \(K\) ; the group \(N\) is isomorphic to the group defined on the set \(K^* \times K\) by the law of composition \((\lambda, \mu)(\lambda', \mu') = (\lambda\lambda', \lambda' \mu + \mu')\) .
\item
  Prove the polynomial identity in 8 indeterminates
\end{enumerate}

\[
\begin{array}{r l} & (\mathrm{X} _ {1} ^ {2} + \mathrm{X} _ {2} ^ {2} + \mathrm{X} _ {3} ^ {2} + \mathrm{X} _ {4} ^ {2}) (\mathrm{Y} _ {1} ^ {2} + \mathrm{Y} _ {2} ^ {2} + \mathrm{Y} _ {3} ^ {2} + \mathrm{Y} _ {4} ^ {2}) \\ & \qquad = (\mathrm{X} _ {1} \mathrm{Y} _ {1} - \mathrm{X} _ {2} \mathrm{Y} _ {2} - \mathrm{X} _ {3} \mathrm{Y} _ {3} - \mathrm{X} _ {4} \mathrm{Y} _ {4}) ^ {2} \\ & \qquad + (\mathrm{X} _ {1} \mathrm{Y} _ {2} + \mathrm{X} _ {2} \mathrm{Y} _ {1} + \mathrm{X} _ {3} \mathrm{Y} _ {4} - \mathrm{X} _ {4} \mathrm{Y} _ {3}) ^ {2} \\ & \qquad + (\mathrm{X} _ {1} \mathrm{Y} _ {3} + \mathrm{X} _ {3} \mathrm{Y} _ {1} + \mathrm{X} _ {4} \mathrm{Y} _ {2} - \mathrm{X} _ {2} \mathrm{Y} _ {4}) ^ {2} \\ & \qquad + (\mathrm{X} _ {1} \mathrm{Y} _ {4} + \mathrm{X} _ {4} \mathrm{Y} _ {1} + \mathrm{X} _ {2} \mathrm{Y} _ {3} - \mathrm{X} _ {3} \mathrm{Y} _ {2}) ^ {2}. \end{array}
\]

(Apply formula (32) of no. 5 to a suitable quaternion algebra.)

\begin{enumerate}
\def\labelenumi{\arabic{enumi}.}
\setcounter{enumi}{19}
\tightlist
\item
  Show that there exists no polynomial identity in 6 indeterminates
\end{enumerate}

\[
(\mathrm{X} _ {1} ^ {2} + \mathrm{X} _ {2} ^ {2} + \mathrm{X} _ {3} ^ {2}) (\mathrm{Y} _ {2} ^ {2} + \mathrm{Y} _ {2} ^ {2} + \mathrm{Y} _ {3} ^ {2}) = \mathrm{P} ^ {2} + \mathrm{Q} ^ {2} + \mathrm{R} ^ {2}
\]

where P, Q, R are three polynomials with coefficients in Z with respect to the indeterminates \(X_{i}\) and \(Y_{i}\) . (Observe that 15 = 3.5 cannot be written in the form \(p^{2} + q^{2} + r^{2}\) , where p, q, r are integers.)

\begin{enumerate}
\def\labelenumi{\arabic{enumi}.}
\setcounter{enumi}{20}
\tightlist
\item
  Let S be a multiplicative monoid with identity element e, satisfying condition (D) of no. 10 and further satisfying the two following conditions: (1) the relation st = e implies s = t = e: (2) for all \(s \in S\) , the number n of terms in a finite sequence \((t_{i})_{1 \leq i \leq n}\) of elements distinct from e and such that \(t_{1}t_{2} \ldots t_{n} = s\) is less than a finite number \(v(s)\) depending only on s.
\end{enumerate}

Show that, for an element \(x = \sum_{s \in S} \alpha_s s\) of the total algebra of \(S\) over a ring \(A\) to be invertible, it is necessary and sufficient that \(\alpha_e\) be invertible in \(A\) (reduce it to the case where \(\alpha_e = 1\) and use the identity

\[
e - z ^ {n + 1} = (e - z) (e + z + \dots + z ^ {n})).
\]

\begin{enumerate}
\def\labelenumi{\arabic{enumi}.}
\setcounter{enumi}{21}
\item
  If K is a commutative field, show that in the ring of formal power series \(K[[X_{1}, X_{2}, \ldots, X_{p}]]\) there is only a single maximal ideal, equal to the set of non-invertible elements. On the other hand, show that in the polynomial ring \(K[X_{1}, \ldots, X_{p}]\) for \(p \geqslant 1\) there always exist several distinct maximal ideals.
\item
  Let K be a commutative field; show that there exists no formal power series \(u(\mathbf{X}, \mathbf{Y}) \in \mathbf{K}[[\mathbf{X}, \mathbf{Y}]]\) such that, for an integer \(m > 0\) ,
\end{enumerate}

\[
(\mathrm{X} + \mathrm{Y}) u (\mathrm{X}, \mathrm{Y}) = \mathrm{X} ^ {m} \mathrm{Y} ^ {m}.
\]

\begin{enumerate}
\def\labelenumi{\arabic{enumi}.}
\setcounter{enumi}{23}
\item
  Let K be a commutative field and let k be an integer such that \(k \neq 0\) in K. Show that for every formal power series u of K{[}{[}X{]}{]} with constant term equal to 1, there exists a formal power series \(v \in K[[X]]\) such that \(v^{k} = u\) .
\item
  Let A be a ring, commutative or otherwise, and let \(\sigma\) be an endomorphism of A. On the additive product group \(\mathbf{E} = \mathbf{A}^{\mathbf{N}}\) an internal law of composition is defined by writing \((\alpha_{n})(\beta_{n}) = (\gamma_{n})\) , where \(\gamma_{n} = \sum_{p + q = n}\alpha_{p}\sigma^{p}(\beta_{q})\) (with the convention \(\sigma^0 (\xi) = \xi\) ).
\end{enumerate}

\begin{enumerate}
\def\labelenumi{(\alph{enumi})}
\tightlist
\item
  Show that this law defines, with addition on E, a ring structure on E, admitting as unit element e the sequence \((\alpha_{n})\) where \(\alpha_{0}=1\) and \(\alpha_{n}=0\) for \(n\geqslant1\) . The mapping which associates with every \(\xi\in A\) the element \((\alpha_{n})\in E\) , such that \(\alpha_{0}=\xi,\alpha_{n}=0\) for \(n\geqslant1\) , is an isomorphism of A onto a
\end{enumerate}

subring of E, with which it is identified. Then we write \(\sum_{n=0}^{\infty} \alpha_n X^n\) instead of

\[
u = \sum_ {n = 0} ^ {\infty} \alpha_ {n} X ^ {n} \neq 0.
\]

\((\alpha_{n})\) and \(\mathbf{X}^{p}\beta = \sigma^{p}(\beta)\mathbf{X}^{p}\) for all \(\beta \in \mathbf{A}\) ; if \(u = \sum_{n=0}^{\infty} \alpha_{n}\mathbf{X}^{n} \neq 0\) , the smallest integer \(n\) such that \(\alpha_{n} \neq 0\) is called the order of \(u\) and denoted by \(\omega(u)\) .

\begin{enumerate}
\def\labelenumi{(\alph{enumi})}
\setcounter{enumi}{1}
\tightlist
\item
  If A is a ring with no divisor of 0 and \(\sigma\) is an isomorphism of A onto a subring of A, show that E is a ring with no divisor of 0 and that
\end{enumerate}

\[
\omega (u v) = \omega (u) + \omega (v)
\]

for \(u \neq 0\) and \(v \neq 0\) .

\begin{enumerate}
\def\labelenumi{(\alph{enumi})}
\setcounter{enumi}{2}
\item
  For \(u = \sum_{n=0}^{\infty} \alpha_n X^n\) to be invertible, it is necessary and sufficient that \(\alpha_0\) be invertible in A.
\item
  Suppose that A is a field and \(\sigma\) an automorphism of A. Show that E admits a field of left fractions (I, §9, Exercise 15) and that in this field F every element \(\neq0\) can be written uniquely in the form \(uX^{-h}\) , where u is an element of E of order 0.
\end{enumerate}

¶ *26. (a) Let A and B be two well-ordered subsets of R (which are necessarily countable; cf.~General Topology, IV, § 2, Exercise 1). Show that A + B is well-ordered and that, for all \(c \in \mathbf{A} + \mathbf{B}\) , there exists only a finite number of ordered pairs \((a, b)\) such that \(a \in \mathbf{A}\) , \(b \in \mathbf{B}\) and \(a + b = c\) (to show that a non-empty subset of A + B has a least element, consider its greatest lower bound in R).

\begin{enumerate}
\def\labelenumi{(\alph{enumi})}
\setcounter{enumi}{1}
\tightlist
\item
  Let K be a commutative field. In the vector space \(K^{R}\) , consider the vector subspace E consisting of the elements \((\alpha_{x})\) such that the set (depending on \((\alpha_{x})\) of \(x \in \mathbf{R}\) such that \(\alpha_{x} \neq 0\) is well-ordered. For any two elements \((\alpha_{x})\) , \((\beta_{x})\) of E, we write \((\alpha_{x})(\beta_{x}) = (\gamma_{x})\) , where \(\gamma_{x} = \sum_{y+z=x} \alpha_{y} \beta_{z}\) (a sum which is meaningful by virtue of (a)). Show that this law of composition defines, with addition, a field structure on E; the elements of E are also denoted by \(\sum_{t \in \mathbb{R}} \alpha_{t} X^{t}\) and called formal power series with well-ordered real exponents, with coefficients in K.*
\end{enumerate}

\BourbakiExercisesSection

\begin{enumerate}
\def\labelenumi{\arabic{enumi}.}
\tightlist
\item
  Let \(\Delta\) be a commutative monoid written additively with the identity element denoted by 0, all of whose elements are cancellable. Let E be a graded A-algebra of type \(\Delta\) and suppose that E admits a unit element \(e = \sum_{\alpha \in \Delta} e_{\alpha}\) where \(e_{\alpha} \in \mathbf{E}_{\alpha}\) . Show that necessarily \(e = e_0 \in \mathbf{E}_0\) .
\end{enumerate}
